%% 3-14 ばらつきを1つの数で：標準偏差（第12回・組 V3-7。ページ数は自由・偶数）
%% この章の部品（名前は \bkthro…）
\makeatletter
% 本文の中の1行見出し  \bkthrohead{見出し}
\newcommand\bkthrohead[1]{\par\addvspace{5pt}\noindent{\sffamily\bfseries\color{bkTheme!85!black}▶\ #1}\par\nopagebreak\vspace{1pt}}
% 数直線  \bkthroaxis{y}{1あたりの mm}{左はしの値}{右はしの値}{目もりの間}{数字を書く間}{小数のけた}
\newcommand\bkthroaxis[7]{%
  \pgfmathsetmacro\bkthro@L{(#4-#3)*#2}%
  \draw[bkInk, line width=0.7pt] (-2,#1) -- (\bkthro@L+2,#1);
  \pgfmathtruncatemacro\bkthro@n{round((#4-#3)/#5)}%
  \foreach \k in {0,...,\bkthro@n} {\pgfmathsetmacro\bkthro@x{\k*#5*#2}\draw[bkInk, line width=0.5pt] (\bkthro@x,#1-1) -- (\bkthro@x,#1+1);}
  \pgfmathtruncatemacro\bkthro@m{round((#4-#3)/#6)}%
  \foreach \k in {0,...,\bkthro@m} {\pgfmathsetmacro\bkthro@x{\k*#6*#2}\pgfmathsetmacro\bkthro@v{#3+\k*#6}%
    \draw[bkInk, line width=0.7pt] (\bkthro@x,#1-1.6) -- (\bkthro@x,#1+1.6);
    \node[anchor=north, font=\sffamily\footnotesize, inner sep=0pt] at (\bkthro@x,#1-2.8) {\pgfmathprintnumber[fixed, fixed zerofill, precision=#7, assume math mode=true]{\bkthro@v}};}}
\makeatother

\bkchapter{3}{14}{ばらつきを1つの\R{数}{かず}で：\\\R{標準偏差}{ひょう|じゅん|へん|さ}}{\bklevel{3}}{ばらつきを1つの\R{数}{かず}で：\R{標準偏差}{ひょう|じゅん|へん|さ}}
\bkmichishirube{2-2,2-6,2-8,3-2,3-13}{3-15,3-16,3-17,9-4}

\begin{goal}
\begin{itemize}[label={\color{\bkTT}●}, leftmargin=1.6\zw, itemsep=0pt, topsep=1pt]
\item \bklevel{3}\enspace \R{偏差}{へん|さ}（\R{値}{あたい}－\R{平均}{へい|きん}）を\R{求}{もと}め、\R{偏差}{へん|さ}をたすと0になるわけを\R{説明}{せつ|めい}できる。
\item \bklevel{3}\enspace \R{分散}{ぶん|さん}と\R{標準偏差}{ひょう|じゅん|へん|さ} $\sigma$ を\R{求}{もと}め、ばらつきの\R{大}{おお}きさをくらべられる。
\item \bklevel{3}\enspace 「\R{平均}{へい|きん}$\pm\sigma$ にだいたい7\R{割}{わり}」と、\R{何回}{なん|かい}もはかった\R{平均}{へい|きん}のずれ $\dfrac{\sigma}{\sqrt{n}}$ を\R{説明}{せつ|めい}できる。
\end{itemize}
\end{goal}

\begin{hatena}
\R{工場}{こう|じょう}で、2\R{台}{だい}の\R{機械}{き|かい}AとBが、ペットボトルにお\R{茶}{ちゃ}を\qty{500}{mL}ずつ\R{入}{い}れています。6\R{本}{ぽん}ずつはかると、\R{平均}{へい|きん}はどちらもちょうど\qty{500}{mL}でしたが、Bのほうがちらばっているようです。「Aは、だいたい\R{何}{なん}mLずれる」「Bは、だいたい\R{何}{なん}mLずれる」を、1つの\R{数}{かず}で\R{言}{い}えるかな？

\begin{bkfigure}
\begin{tikzpicture}[x=1mm, y=1mm]
  % 1 mL を 6 mm、492 mL が 0
  \bkthroaxis{0}{6}{492}{508}{1}{2}{0}
  \foreach \v in {501,497,503,500,498} {\fill[bkBlue] ({(\v-492)*6},14) circle (1.3);}
  \fill[bkBlue] ({(501-492)*6},17) circle (1.3);
  \foreach \v in {506,496,500,494} {\fill[bkOrange] ({(\v-492)*6},5) circle (1.3);}
  \fill[bkOrange] ({(502-492)*6},5) circle (1.3) ({(502-492)*6},8) circle (1.3);
  \draw[bkRed!80!black, line width=0.8pt, dash pattern=on 2pt off 1pt] (48,-0.5) -- (48,21);
  \node[anchor=east, font=\sffamily\bfseries\small, inner sep=0pt] at (-5,14) {A};
  \node[anchor=east, font=\sffamily\bfseries\small, inner sep=0pt] at (-5,5) {B};
  \node[anchor=south, font=\sffamily\footnotesize, text=bkRed!80!black, inner sep=0pt] at (48,22) {\R{平均}{へい|きん}500};
  \node[anchor=north, font=\sffamily\footnotesize, inner sep=0pt] at (48,-7.6) {お\R{茶}{ちゃ}の\R{量}{りょう}（mL）。1\R{本}{ぽん}が\R{点}{てん}1こ};
\end{tikzpicture}
\end{bkfigure}
\end{hatena}

%% ---------------- やってみよう・思い出そう ----------------
\begin{yattemiyou}[ずれを\R{全部}{ぜん|ぶ}たしてみよう]
\R{機械}{き|かい}Aの6\R{本}{ぽん}の\R{量}{りょう}（mL）です。
\begin{bkcenter}
\begin{bktable}{colspec={Q[l,m,wd=30mm] *{6}{Q[c,m,wd=12mm]}}, rowsep=1pt, colsep=2pt, cells={font=\footnotesize}, row{1}={font=\sffamily\footnotesize, bg=bkGreen!12}}
\R{何本}{なん|ぼん}め & 1 & 2 & 3 & 4 & 5 & 6 \\
\R{量}{りょう}（mL） & 501 & 497 & 503 & 500 & 498 & 501 \\
\R{平均}{へい|きん}とのちがい & & & & & & \\
\end{bktable}
\end{bkcenter}
\begin{enumerate}[label={\bkcircled[bkGreen]{\arabic*}}, leftmargin=2.4\zw, itemsep=1pt, topsep=2pt]
\item 6\R{本}{ぽん}の\R{平均}{へい|きん}が\qty{500}{mL}になることを、\R{計算}{けい|さん}で\R{確}{たし}かめよう。
\item それぞれの\R{量}{りょう}から\R{平均}{へい|きん}500をひいて、「\R{平均}{へい|きん}とのちがい」をノートに\R{書}{か}こう。\R{平均}{へい|きん}より\R{少}{すく}ないときは、マイナスの\R{数}{かず}になるよ（2-2）。
\item 「\R{平均}{へい|きん}とのちがい」を6こ\R{全部}{ぜん|ぶ}たしてみよう。
\item \R{機械}{き|かい}Bの6\R{本}{ぽん}（502、506、496、500、494、502）でも、\R{同}{おな}じことをしよう。
\end{enumerate}
\end{yattemiyou}

\begin{omoidasou}
\bkrecall{3-13}{\R{平均}{へい|きん} $\bar{x}=\dfrac{x_1+x_2+\cdots+x_n}{n}$。\R{範囲}{はん|い}＝\R{最大値}{さい|だい|ち}－\R{最小値}{さい|しょう|ち}}
\bkrecall{2-2}{マイナス×マイナスはプラス。$(-3)^2=(-3)\times(-3)=9$}
\bkrecall{2-8}{$\sqrt{\text{□}}$ は、2\R{乗}{じょう}すると□になる\R{正}{せい}の\R{数}{かず}。$\sqrt{16}=4$、$\sqrt{2}=1.414\cdots$}
\bkrecall{3-2}{\R{絶対値}{ぜっ|たい|ち}は0からの\R{距離}{きょ|り}。$|-3|=3$、$|3|=3$}
\bkrecall{2-6}{「\qty{5.0}{cm}」は $\pm$\qty{0.05}{cm}のずれをふくむ。$\pm$ はずれのはば。}
\end{omoidasou}

\bkfukidashi{sora}{あれ？\ AもBも、たすと0になっちゃった。Bのほうがずれているはずなのに……。}

\clearpage
\bkthrohead{この\R{章}{しょう}でやること}
3-13では、データのちらばりを、ヒストグラムや\R{箱}{はこ}ひげ\R{図}{ず}で\R{目}{め}に\R{見}{み}えるようにしました。この\R{章}{しょう}では、「ちらばりの\R{大}{おお}きさ」を1つの\R{数}{かず}で\R{表}{あらわ}します。この\R{数}{かず}があると、「この\R{機械}{き|かい}は、だいたい\qty{2}{mL}ずれる」「このはかり\R{方}{かた}は、だいたい0.1\R{秒}{びょう}ずれる」と、ずれの\R{大}{おお}きさを\R{言}{い}えるようになります。

%% ---------------- 数学の道具 ----------------
\bksection{doguu}{\R{数学}{すう|がく}の\R{道具}{どう|ぐ}}

\bksubsection{1}{\R{平均}{へい|きん}からのずれ：\R{偏差}{へん|さ}}

\leavevmode\label{bk:3-14:hensa}やってみようの②で\R{求}{もと}めた「\R{平均}{へい|きん}とのちがい」を、\textbf{\R{偏差}{へん|さ}}といいます。\bkidx{へんさ}{\R{偏差}{へん|さ}}「\R{偏}{へん}」は「かたよる」、「\R{差}{さ}」は「ひき\R{算}{ざん}の\R{答}{こた}え」という\R{意味}{い|み}です。
\begin{center}
\R{偏差}{へん|さ}＝\R{値}{あたい}－\R{平均}{へい|きん}\qquad（\R{文字}{も|じ}で\R{書}{か}くと $x_1-\bar{x}$、$x_2-\bar{x}$、……）
\end{center}
\R{平均}{へい|きん}より\R{大}{おお}きい\R{値}{あたい}の\R{偏差}{へん|さ}はプラス、\R{小}{ちい}さい\R{値}{あたい}の\R{偏差}{へん|さ}はマイナス、\R{平均}{へい|きん}と\R{同}{おな}じ\R{値}{あたい}の\R{偏差}{へん|さ}は0です。\R{機械}{き|かい}Aで\R{計算}{けい|さん}すると、\R{次}{つぎ}のようになります。

\begin{bkcenter}
\begin{bktable}{colspec={Q[l,m,wd=30mm] *{6}{Q[c,m,wd=12mm]} Q[c,m,wd=14mm]}, rowsep=1pt, colsep=2pt, cells={font=\footnotesize}, row{1}={font=\sffamily\footnotesize, bg=bkTheme!12}, column{8}={bg=bkGray!10}}
\R{何本}{なん|ぼん}め & 1 & 2 & 3 & 4 & 5 & 6 & \R{合計}{ごう|けい} \\
\R{量}{りょう}（mL） & 501 & 497 & 503 & 500 & 498 & 501 & 3000 \\
\R{偏差}{へん|さ}（mL） & $+1$ & $-3$ & $+3$ & 0 & $-2$ & $+1$ & 0 \\
\end{bktable}
\bktabcaption{\R{機械}{き|かい}Aの\R{偏差}{へん|さ}。$501-500=+1$、$497-500=-3$、……}
\end{bkcenter}

\R{偏差}{へん|さ}を\R{数直線}{すう|ちょく|せん}に\R{矢印}{や|じるし}でかくと、\R{平均}{へい|きん}から\R{右}{みぎ}へ\R{出}{で}る\R{矢印}{や|じるし}と、\R{左}{ひだり}へ\R{出}{で}る\R{矢印}{や|じるし}があります。

\begin{bkfigure}
\begin{tikzpicture}[x=1mm, y=1mm]
  % 1 mL を 10 mm、495 が 0
  \bkthroaxis{0}{10}{495}{505}{1}{1}{0}
  \draw[bkRed!80!black, line width=0.8pt, dash pattern=on 2pt off 1pt] (50,-0.5) -- (50,44);
  \node[anchor=south, font=\sffamily\footnotesize, text=bkRed!80!black, inner sep=0pt] at (50,45) {\R{平均}{へい|きん}500};
  \foreach \i/\v/\t in {1/501/{+1},2/497/{-3},3/503/{+3},4/500/{0},5/498/{-2},6/501/{+1}} {
    {\pgfmathsetmacro\yy{42-\i*6}
     \node[anchor=east, font=\sffamily\scriptsize, text=bkGray!90!black, inner sep=0pt] at (-4,\yy) {\i\R{本}{ほん}め};
     \ifnum\v=500 \fill[bkInk] (50,\yy) circle (1.1);
     \else\ifnum\v>500 \draw[bkBlue, line width=1.1pt, -{Stealth[length=3pt]}] (50,\yy) -- ({(\v-495)*10},\yy);
     \else \draw[bkOrange!90!black, line width=1.1pt, -{Stealth[length=3pt]}] (50,\yy) -- ({(\v-495)*10},\yy);\fi\fi
     \node[anchor=west, font=\sffamily\footnotesize, inner sep=0pt] at (104,\yy) {$\t$};}}
  \node[anchor=south west, font=\sffamily\footnotesize, inner sep=0pt] at (100,40) {\R{偏差}{へん|さ}};
\end{tikzpicture}
\bkfigcaption{\R{青}{あお}は\R{平均}{へい|きん}より\R{大}{おお}きい（プラス）、だいだいは\R{小}{ちい}さい（マイナス）。4\R{本}{ほん}めは\R{平均}{へい|きん}ちょうどで0。}
\end{bkfigure}

\bkthrohead{\R{偏差}{へん|さ}をたすと、いつも0になる}
やってみようの③で、\R{偏差}{へん|さ}を\R{全部}{ぜん|ぶ}たすと、AもBも0になりました。たまたまではありません。どんなデータでも、\R{偏差}{へん|さ}の\R{合計}{ごう|けい}はかならず0になります。そのわけを、2\R{通}{とお}りで\R{見}{み}てみましょう。

\noindent\textbf{わけ1：つみ\R{木}{き}で\R{見}{み}る}\quad 2-6では、\R{平均}{へい|きん}を「でこぼこをならした\R{高}{たか}さ」と\R{考}{かんが}えました。\R{偏差}{へん|さ}を\R{棒}{ぼう}にすると、\R{平均}{へい|きん}より\R{上}{うえ}の\R{出}{で}っぱりと、\R{下}{した}のへこみに\R{分}{わ}かれます。

\begin{bkfigure}
\begin{tikzpicture}[x=1mm, y=1mm]
  % 偏差1を5mm
  \draw[bkRed!80!black, line width=0.8pt, dash pattern=on 2pt off 1pt] (-4,0) -- (82,0);
  \node[anchor=east, font=\sffamily\footnotesize, text=bkRed!80!black, inner sep=0pt] at (-5,0) {\R{平均}{へい|きん}};
  \foreach \i/\d in {1/1,2/-3,3/3,4/0,5/-2,6/1} {
    {\pgfmathsetmacro\xx{\i*12}
     \ifnum\d>0 \filldraw[fill=bkBlue!30, draw=bkBlue!80!black, line width=0.6pt] (\xx-4,0) rectangle (\xx+4,{\d*5});
       \node[anchor=south, font=\sffamily\footnotesize, inner sep=0pt] at (\xx,{\d*5+1.2}) {$+\d$};
     \else\ifnum\d<0 \filldraw[fill=bkOrange!35, draw=bkOrange!80!black, line width=0.6pt] (\xx-4,0) rectangle (\xx+4,{\d*5});
       \node[anchor=north, font=\sffamily\footnotesize, inner sep=0pt] at (\xx,{\d*5-1.2}) {$\d$};
     \else \node[anchor=south, font=\sffamily\footnotesize, inner sep=0pt] at (\xx,1.2) {0};\fi\fi}}
  \node[anchor=west, font=\sffamily\footnotesize, text=bkBlue!80!black, align=left, inner sep=0pt] at (88,9) {\R{出}{で}っぱりの\R{合計}{ごう|けい}\\$1+3+1=5$};
  \node[anchor=west, font=\sffamily\footnotesize, text=bkOrange!80!black, align=left, inner sep=0pt] at (88,-9) {へこみの\R{合計}{ごう|けい}\\$3+2=5$};
\end{tikzpicture}
\bkfigcaption{\R{平均}{へい|きん}の\R{線}{せん}より\R{上}{うえ}の\R{出}{で}っぱりと、\R{下}{した}のへこみは、\R{同}{おな}じ\R{大}{おお}きさ}
\end{bkfigure}

\R{出}{で}っぱりを\R{切}{き}り\R{取}{と}ってへこみをうめると、ちょうどたいらになります。それが\R{平均}{へい|きん}でした。だから、\R{出}{で}っぱりの\R{合計}{ごう|けい}（プラスの\R{偏差}{へん|さ}の\R{合計}{ごう|けい}）と、へこみの\R{合計}{ごう|けい}（マイナスの\R{偏差}{へん|さ}の\R{合計}{ごう|けい}）は\R{同}{おな}じ\R{大}{おお}きさです。たすと、$5+(-5)=0$ です。

\noindent\textbf{わけ2：\R{式}{しき}で\R{見}{み}る}\quad \R{偏差}{へん|さ}を\R{全部}{ぜん|ぶ}たした\R{式}{しき}を、\R{値}{あたい}の\R{部分}{ぶ|ぶん}と $\bar{x}$ の\R{部分}{ぶ|ぶん}に\R{分}{わ}けて\R{並}{なら}べかえます（たし\R{算}{ざん}とひき\R{算}{ざん}は、\R{順番}{じゅん|ばん}を\R{入}{い}れかえてよい。2-12）。
\begin{align*}
&(x_1-\bar{x})+(x_2-\bar{x})+\cdots+(x_n-\bar{x})\\
&=(x_1+x_2+\cdots+x_n)-(\bar{x}+\bar{x}+\cdots+\bar{x})
\end{align*}
$\bar{x}$ は $n$ こあるので、うしろのかっこは $n\bar{x}$ です。また、\R{平均}{へい|きん}の\R{式}{しき}
\[ \bar{x}=\frac{x_1+x_2+\cdots+x_n}{n} \]
の\R{両辺}{りょう|へん}に $n$ をかけると（2-14）、$x_1+x_2+\cdots+x_n=n\bar{x}$ です。だから、
\[ (x_1+x_2+\cdots+x_n)-n\bar{x}=n\bar{x}-n\bar{x}=0 \]
\R{機械}{き|かい}Aで\R{確}{たし}かめると、$3000-6\times500=3000-3000=0$。\R{式}{しき}のとおりです。

\noindent\textbf{だから、どうなる？}\quad \R{偏差}{へん|さ}をそのままたすと、ばらつきが\R{大}{おお}きくても\R{小}{ちい}さくても、いつも0です。これでは、ばらつきの\R{大}{おお}きさを\R{表}{あらわ}せません。プラスとマイナスが\R{打}{う}ち\R{消}{け}しあってしまうのです。

\bksubsection{2}{2\R{乗}{じょう}して\R{平均}{へい|きん}し、$\sqrt{\phantom{2}}$ でもどす：\R{標準偏差}{ひょう|じゅん|へん|さ}}

\leavevmode\label{bk:3-14:sigma}\R{打}{う}ち\R{消}{け}しあわないようにするには、マイナスをなくせばよいのです。その\R{方法}{ほう|ほう}は2つあります。

\bkthrohead{\R{方法}{ほう|ほう}1：\R{絶対値}{ぜっ|たい|ち}にしてから\R{平均}{へい|きん}する}
\R{絶対値}{ぜっ|たい|ち}は0からの\R{距離}{きょ|り}なので、マイナスになりません（3-2）。$|-3|=3$ です。\R{偏差}{へん|さ}の\R{絶対値}{ぜっ|たい|ち}は、「\R{平均}{へい|きん}から\R{何}{なん}mLはなれているか」そのものです。
\begin{align*}
\text{A}\quad&(1+3+3+0+2+1)\div6=10\div6=1.66\cdots\\
\text{B}\quad&(2+6+4+0+6+2)\div6=20\div6=3.33\cdots
\end{align*}
Bのほうが2\R{倍}{ばい}ばらついている、とわかります。

\bkthrohead{\R{方法}{ほう|ほう}2：2\R{乗}{じょう}してから\R{平均}{へい|きん}する}
マイナスの\R{数}{かず}も、2\R{乗}{じょう}するとプラスになります（2-2）。$(-3)^2=9$ です。\R{機械}{き|かい}Aの\R{偏差}{へん|さ}を2\R{乗}{じょう}して、\R{平均}{へい|きん}します。

\begin{bkcenter}
\begin{bktable}{colspec={Q[l,m,wd=30mm] *{6}{Q[c,m,wd=12mm]} Q[c,m,wd=14mm]}, rowsep=1pt, colsep=2pt, cells={font=\footnotesize}, row{1}={font=\sffamily\footnotesize, bg=bkTheme!12}, column{8}={bg=bkGray!10}}
\R{何本}{なん|ぼん}め & 1 & 2 & 3 & 4 & 5 & 6 & \R{合計}{ごう|けい} \\
\R{偏差}{へん|さ} & $+1$ & $-3$ & $+3$ & 0 & $-2$ & $+1$ & 0 \\
\R{偏差}{へん|さ}の2\R{乗}{じょう} & 1 & 9 & 9 & 0 & 4 & 1 & 24 \\
\end{bktable}
\bktabcaption{$(-3)^2=(-3)\times(-3)=9$、$(-2)^2=4$。2\R{乗}{じょう}すると、どれもプラスか0}
\end{bkcenter}
\[ (1+9+9+0+4+1)\div6=24\div6=4 \]

この「\R{偏差}{へん|さ}の2\R{乗}{じょう}の\R{平均}{へい|きん}」を、\textbf{\R{分散}{ぶん|さん}}といいます。\bkidx{ぶんさん}{\R{分散}{ぶん|さん}}\R{機械}{き|かい}Aの\R{分散}{ぶん|さん}は4です。

\bkthrohead{$\sqrt{\phantom{2}}$ で、もとの\R{単位}{たん|い}にもどす}
\R{分散}{ぶん|さん}の\R{単位}{たん|い}は\R{何}{なん}でしょう。\R{偏差}{へん|さ}はmLなので、2\R{乗}{じょう}すると「mL×mL」です。cm×cmが\R{面積}{めん|せき}の\R{単位}{たん|い}\unit{cm^2}になるのと\R{同}{おな}じで、もとのmLとはちがう\R{単位}{たん|い}になってしまいました。\R{分散}{ぶん|さん}の4は「\qty{4}{mL}ずれる」という\R{意味}{い|み}ではありません。

そこで、2\R{乗}{じょう}したのをもとにもどすために、$\sqrt{\phantom{2}}$ をとります（2-8：2\R{乗}{じょう}すると4になる\R{正}{せい}の\R{数}{かず}）。
\[ \sqrt{4}=2 \]
これで\R{単位}{たん|い}がmLにもどり、「\R{機械}{き|かい}Aは、だいたい\qty{2}{mL}ずれる」と\R{言}{い}えます。この\R{数}{かず}を\textbf{\R{標準偏差}{ひょう|じゅん|へん|さ}}といい、ギリシャ\R{文字}{も|じ}の $\sigma$（シグマ）で\R{表}{あらわ}します。\bkidx{ひょうじゅんへんさ}{\R{標準偏差}{ひょう|じゅん|へん|さ}}「\R{標準}{ひょう|じゅん}」は「ふつうの、\R{目安}{め|やす}になる」という\R{意味}{い|み}です。

\bkthrohead{\R{図}{ず}で\R{見}{み}る：\R{正方形}{せい|ほう|けい}をならす}
\R{偏差}{へん|さ}の2\R{乗}{じょう}は、1\R{辺}{ぺん}が\R{偏差}{へん|さ}の\R{大}{おお}きさの\R{正方形}{せい|ほう|けい}の\R{面積}{めん|せき}です。\R{分散}{ぶん|さん}は、その\R{面積}{めん|せき}をならした\R{平均}{へい|きん}です。$\sigma$ は、\R{面積}{めん|せき}が\R{分散}{ぶん|さん}になる\R{正方形}{せい|ほう|けい}の1\R{辺}{ぺん}です。

\begin{bkfigure}
\begin{tikzpicture}[x=1mm, y=1mm]
  % 1 mL を 4 mm
  \foreach \x/\d/\l in {0/1/{+1},8/3/{-3},24/3/{+3},40/0/{0},48/2/{-2},60/1/{+1}} {
    {\ifnum\d>0 \filldraw[fill=bkBlue!25, draw=bkBlue!80!black, line width=0.6pt] (\x,0) rectangle (\x+\d*4,\d*4);
       \node[anchor=center, font=\sffamily\scriptsize, inner sep=0pt] at (\x+\d*2,\d*2) {\pgfmathparse{int(\d*\d)}\pgfmathresult};
     \else \fill[bkInk] (\x+1,0) circle (0.6);\fi
     \node[anchor=north, font=\sffamily\scriptsize, inner sep=0pt] at (\x+\d*2,-1.4) {$\l$};}}
  \node[anchor=north, font=\sffamily\footnotesize, inner sep=0pt] at (34,-6.4) {6つの\R{正方形}{せい|ほう|けい}（\R{面積}{めん|せき}の\R{合計}{ごう|けい}24）};
  \draw[bk arrow, bkGray, line width=0.9pt] (70,6) -- (82,6);
  \node[anchor=south, font=\sffamily\footnotesize, align=center, inner sep=0pt] at (76,8) {ならす};
  \filldraw[fill=bkGreen!30, draw=bkGreen!70!black, line width=0.8pt] (88,0) rectangle (96,8);
  \node[font=\sffamily\scriptsize, inner sep=0pt] at (92,4) {4};
  \draw[bkGreen!60!black, line width=0.6pt, decorate, decoration={brace, amplitude=1mm, mirror}] (88,-1) -- (96,-1);
  \node[anchor=north, font=\sffamily\footnotesize, text=bkGreen!50!black, inner sep=0pt] at (92,-3.4) {1\R{辺}{ぺん}2＝$\sigma$};
  \node[anchor=west, font=\sffamily\footnotesize, align=left, inner sep=0pt] at (100,5) {$24\div6=4$\\（\R{分散}{ぶん|さん}）};
\end{tikzpicture}
\bkfigcaption{\R{正方形}{せい|ほう|けい}の\R{中}{なか}の\R{数}{かず}は\R{面積}{めん|せき}（\R{偏差}{へん|さ}の2\R{乗}{じょう}）。ならした\R{正方形}{せい|ほう|けい}の\R{面積}{めん|せき}が\R{分散}{ぶん|さん}、1\R{辺}{ぺん}が $\sigma$。}
\end{bkfigure}

\bkthrohead{\R{機械}{き|かい}Bでも\R{計算}{けい|さん}する}
Bの\R{偏差}{へん|さ}は $+2$、$+6$、$-4$、0、$-6$、$+2$ です。2\R{乗}{じょう}すると4、36、16、0、36、4。
\[ \text{\R{分散}{ぶん|さん}}=(4+36+16+0+36+4)\div6=96\div6=16,\qquad \sigma=\sqrt{16}=4 \]
\bkhatenaref の\R{答}{こた}え：\R{機械}{き|かい}Aは $\sigma=2$ で「だいたい\qty{2}{mL}ずれる」、\R{機械}{き|かい}Bは $\sigma=4$ で「だいたい\qty{4}{mL}ずれる」。Bのばらつきは、Aの2\R{倍}{ばい}です。

\begin{taisetsu}
\R{標準偏差}{ひょう|じゅん|へん|さ} $\sigma$ の\R{求}{もと}め\R{方}{かた}\\
① \R{平均}{へい|きん} $\bar{x}$ を\R{求}{もと}める\quad ② \R{偏差}{へん|さ}（\R{値}{あたい}－\R{平均}{へい|きん}）を\R{求}{もと}める\quad ③ \R{偏差}{へん|さ}を2\R{乗}{じょう}する\\
④ 2\R{乗}{じょう}の\R{平均}{へい|きん}を\R{求}{もと}める（これが\R{分散}{ぶん|さん}）\quad ⑤ \R{分散}{ぶん|さん}の $\sqrt{\phantom{2}}$ をとる（これが $\sigma$）
\end{taisetsu}

①〜⑤を1つの\R{式}{しき}にまとめると、\R{次}{つぎ}のようになります。\R{覚}{おぼ}えなくても、①〜⑤の\R{順}{じゅん}に\R{計算}{けい|さん}すれば\R{求}{もと}められます。
\[ \sigma=\sqrt{\frac{(x_1-\bar{x})^2+(x_2-\bar{x})^2+\cdots+(x_n-\bar{x})^2}{n}} \]

\noindent\begin{minipage}[c]{0.47\linewidth}
$\sigma$ は、ギリシャ\R{文字}{も|じ}のシグマの\R{小文字}{こ|も|じ}です。\R{同}{おな}じシグマでも、\R{大文字}{おお|も|じ}は\R{第}{だい}5\R{巻}{かん}（5-10）で「たし\R{算}{ざん}をまとめて\R{書}{か}く\R{記号}{き|ごう}」として\R{出}{で}てくる、\R{別}{べつ}の\R{意味}{い|み}の\R{記号}{き|ごう}です。まぜないようにしましょう。
\end{minipage}\hfill
\begin{minipage}[c]{0.5\linewidth}\raggedleft
\bkkigocard{$\sigma$}{シグマ}{\R{標準偏差}{ひょう|じゅん|へん|さ}。\R{偏差}{へん|さ}の2\R{乗}{じょう}の\R{平均}{へい|きん}（\R{分散}{ぶん|さん}）の $\sqrt{\phantom{2}}$。ばらつきの\R{大}{おお}きさ}{sigma（standard deviation）}{3-14}
\end{minipage}

\bkthrohead{\R{表計算}{ひょう|けい|さん}で $\sigma$ を\R{求}{もと}める}
3-13のふりこの20\R{回}{かい}（\R{平均}{へい|きん}14.20\R{秒}{びょう}）でも、①〜⑤をしてみます。\R{数}{かず}が\R{多}{おお}いので、\R{表計算}{ひょう|けい|さん}（2-18）を\R{使}{つか}います。

\begin{bkcenter}
\bksheet{4}{24}{%
1 & \R{番号}{ばん|ごう} & \R{時間}{じ|かん}（\R{秒}{びょう}） & \R{偏差}{へん|さ} & \R{偏差}{へん|さ}の2\R{乗}{じょう} \\
2 & 1 & 14.37 & 0.17 & 0.0289 \\
3 & 2 & 14.29 & 0.09 & 0.0081 \\
4 & 3 & 14.24 & 0.04 & 0.0016 \\
5 & 4 & 14.12 & $-0.08$ & 0.0064 \\
$\vdots$ & $\vdots$ & $\vdots$ & $\vdots$ & $\vdots$ \\
21 & 20 & 14.05 & $-0.15$ & 0.0225 \\
22 & & \R{合計}{ごう|けい} & 0.00 & 0.2348 \\
23 & & & \R{分散}{ぶん|さん} & 0.01174 \\
24 & & & $\sigma$ & 0.108 \\}
\bktabcaption{C2に \texttt{=B2-14.20}（\R{偏差}{へん|さ}）、D2に \texttt{=C2\^{}2}（2\R{乗}{じょう}）と\R{入}{い}れて\R{下}{した}へコピー。D23は \texttt{=D22/20}、D24は \texttt{=SQRT(D23)}}
\end{bkcenter}

\texttt{SQRT} は $\sqrt{\phantom{2}}$ のことです。ふりこの $\sigma$ は0.108\R{秒}{びょう}、およそ0.11\R{秒}{びょう}です。\texttt{=STDEV.P(B2:B21)} と\R{入}{い}れると、①〜⑤を1\R{回}{かい}で\R{計算}{けい|さん}してくれます（P については、この\R{章}{しょう}の\R{最後}{さい|ご}の「せいかくに\R{言}{い}うと」を\R{見}{み}てください）。

\bkthrohead{なぜ\R{絶対値}{ぜっ|たい|ち}ではなく、2\R{乗}{じょう}を\R{使}{つか}うの？}
\R{方法}{ほう|ほう}1の\R{絶対値}{ぜっ|たい|ち}でも、ばらつきはくらべられました。それでも、\R{理科}{り|か}や\R{工学}{こう|がく}では、ふつう2\R{乗}{じょう}を\R{使}{つか}う $\sigma$ が\R{使}{つか}われます。
\begin{enumerate}[label={\bkcircled[bkBlue]{\arabic*}}, leftmargin=2.4\zw, itemsep=1pt, topsep=2pt]
\item \textbf{\R{大}{おお}きなずれを\R{重}{おも}く\R{見}{み}る。}\quad ずれが2\R{倍}{ばい}になると、2\R{乗}{じょう}は4\R{倍}{ばい}になります。\R{部品}{ぶ|ひん}づくりでは、\R{小}{ちい}さなずれがたくさんあるより、とても\R{大}{おお}きなずれが1つあるほうがこまるので、それを\R{大}{おお}きく\R{数}{かぞ}えるのは\R{理}{り}にかなっています。
\item \textbf{きまりがきれいに\R{成}{な}り\R{立}{た}つ。}\quad このあとの④⑤の「\R{単位}{たん|い}を\R{変}{か}えたときのきまり」「\R{何回}{なん|かい}もはかった\R{平均}{へい|きん}のずれは $\dfrac{\sigma}{\sqrt{n}}$」は、2\R{乗}{じょう}を\R{使}{つか}ったからこそ\R{計算}{けい|さん}で\R{作}{つく}れます。3-16の「ずれがいちばん\R{小}{ちい}さくなる\R{線}{せん}」も、2\R{乗}{じょう}を\R{使}{つか}うので\R{平方完成}{へい|ほう|かん|せい}（3-8）で\R{求}{もと}められます。
\item \textbf{\R{場合}{ば|あい}\R{分}{わ}けがいらない。}\quad \R{絶対値}{ぜっ|たい|ち}は「プラスならそのまま、マイナスならマイナスの\R{記号}{き|ごう}をとる」と\R{分}{わ}けて\R{考}{かんが}えますが、2\R{乗}{じょう}はいつも\R{同}{おな}じ\R{計算}{けい|さん}です。
\end{enumerate}

\bksubsection{3}{\R{平均}{へい|きん}$\pm\sigma$ の\R{中}{なか}に、だいたい7\R{割}{わり}}

\leavevmode\label{bk:3-14:nanawari}$\sigma$ は「だいたいこのくらいずれる」という\R{数}{かず}でした。では、「だいたい」とは、どのくらいでしょう。ふりこの20\R{回}{かい}で\R{調}{しら}べます。\R{平均}{へい|きん}は14.20\R{秒}{びょう}、$\sigma$ はおよそ0.11\R{秒}{びょう}でした。

\R{平均}{へい|きん}から $\sigma$ だけ\R{下}{した}と\R{上}{うえ}は、
\[ 14.20-0.11=14.09,\qquad 14.20+0.11=14.31 \]
です。この\R{範囲}{はん|い}を「\R{平均}{へい|きん}$\pm\sigma$」と\R{書}{か}きます。20\R{回}{かい}の\R{点}{てん}をかいて、この\R{範囲}{はん|い}に\R{入}{はい}る\R{点}{てん}を\R{数}{かぞ}えましょう。

\begin{bkfigure}
\begin{tikzpicture}[x=1mm, y=1mm]
  % 0.01秒を2mm、13.90 が 0
  \fill[bkGreen!12] ({(13.98-13.9)*200},-0.6) rectangle ({(14.42-13.9)*200},21.5);
  \fill[bkGreen!30] ({(14.09-13.9)*200},-0.6) rectangle ({(14.31-13.9)*200},14);
  \bkthroaxis{0}{200}{13.90}{14.50}{0.05}{0.1}{2}
  \foreach \v in {14.37,14.29,14.24,14.12,14.21,14.17,14.33,14.15,14.22,14.19,13.95,14.18,14.23,14.27,14.07,14.20,14.44,14.05} {
    \pgfmathsetmacro\c{(\v>=14.09 && \v<=14.31) ? 1 : 0}
    \ifdim\c pt>0.5pt \fill[bkBlue] ({(\v-13.9)*200},2.6) circle (1); \else \fill[bkRed] ({(\v-13.9)*200},2.6) circle (1); \fi}
  \fill[bkBlue] ({(14.16-13.9)*200},2.6) circle (1) ({(14.16-13.9)*200},5.0) circle (1);
  \node[anchor=south, font=\sffamily\footnotesize, text=bkGreen!30!black, inner sep=0pt] at (60,8.6) {\R{平均}{へい|きん}$\pm\sigma$：14こ};
  \node[anchor=south, font=\sffamily\footnotesize, text=bkGreen!30!black, inner sep=0pt] at (60,16.6) {\R{平均}{へい|きん}$\pm2\sigma$：18こ};
  \node[anchor=west, font=\sffamily\footnotesize, inner sep=0pt] at (52,-11) {\R{時間}{じ|かん}（\R{秒}{びょう}）};
\end{tikzpicture}
\bkfigcaption{\R{青}{あお}は\R{平均}{へい|きん}$\pm\sigma$ の\R{中}{なか}、\R{赤}{あか}は\R{外}{そと}。こい\R{緑}{みどり}が\R{平均}{へい|きん}$\pm\sigma$、うすい\R{緑}{みどり}が\R{平均}{へい|きん}$\pm2\sigma$。}
\end{bkfigure}

20このうち14こが、\R{平均}{へい|きん}$\pm\sigma$ の\R{中}{なか}に\R{入}{はい}っています。$14\div20=0.7$ で、7\R{割}{わり}です。\R{外}{そと}にあるのは6こ（14.37、14.33、14.07、14.05、13.95、14.44）で、3\R{割}{わり}です。

ヒストグラムがつりがね\R{形}{がた}（まん\R{中}{なか}が\R{高}{たか}く、\R{左右}{さ|ゆう}に\R{同}{おな}じように\R{低}{ひく}くなる\R{山}{やま}）のデータでは、いつも、
\begin{center}
\R{平均}{へい|きん}$\pm\sigma$ の\R{中}{なか}に、だいたい7\R{割}{わり}（くわしくは68\%）
\end{center}
が\R{入}{はい}ることが\R{知}{し}られています。$\pm2\sigma$ まで\R{広}{ひろ}げると、だいたい95\%が\R{入}{はい}ります（ふりこでは18こで90\%。\R{数}{かず}を\R{多}{おお}くすると95\%に\R{近}{ちか}づきます）。なぜ68\%や95\%になるのかは、3-17の「\R{正規分布}{せい|き|ぶん|ぷ}」で\R{学}{まな}びます。

\noindent\textbf{つりがね\R{形}{がた}でないときは？}\quad 3-13のものさしキャッチのように、1つだけ\R{大}{おお}きくはなれた\R{値}{あたい}があるデータや、\R{山}{やま}が\R{左右}{さ|ゆう}どちらかにかたよったデータでは、7\R{割}{わり}になるとはかぎりません。「だいたい7\R{割}{わり}」は、つりがね\R{形}{がた}のときの\R{目安}{め|やす}です。

\bkthrohead{2-6の「$\pm$」とのちがい}
2-6では、「\qty{5.0}{cm}」は $\pm$\qty{0.05}{cm}のずれをふくむ、と\R{学}{まな}びました。\R{同}{おな}じ「$\pm$」でも、\R{意味}{い|み}がちがいます。

\begin{bkcenter}
\begin{bktable}{colspec={Q[l,m,wd=26mm] Q[l,m,wd=48mm] Q[l,m,wd=48mm]}, rowsep=1.2pt, colsep=3pt, cells={font=\footnotesize}, row{1}={font=\sffamily\bfseries\footnotesize, bg=bkTheme!12}}
 & 2-6の $\pm$ & この\R{章}{しょう}の $\pm\sigma$ \\
\R{例}{れい} & \qty{5.0}{cm}（$\pm$\qty{0.05}{cm}） & 14.20\R{秒}{びょう}$\pm$0.11\R{秒}{びょう}（$\pm\sigma$） \\
\R{何}{なに}を\R{表}{あらわ}す？ & \R{目}{め}もりを\R{読}{よ}むときの、いちばん\R{大}{おお}きいずれ & たくさんはかったときの、ふつうのずれ \\
\R{中}{なか}に\R{入}{はい}るのは？ & \R{全部}{ぜん|ぶ}（はみ\R{出}{だ}さない） & だいたい7\R{割}{わり}（3\R{割}{わり}ははみ\R{出}{だ}す） \\
\end{bktable}
\bktabcaption{\R{同}{おな}じ「$\pm$」でも、どちらの\R{意味}{い|み}かを\R{書}{か}きそえる}
\end{bkcenter}

\R{実験}{じっ|けん}のレポートや\R{工業製品}{こう|ぎょう|せい|ひん}の\R{説明書}{せつ|めい|しょ}では、「$\pm$」のうしろの\R{数}{かず}が\R{何}{なに}を\R{表}{あらわ}すかを\R{書}{か}きそえます。

\bksubsection{4}{\R{単位}{たん|い}を\R{変}{か}えると、$\sigma$ はどうなる？}

\leavevmode\label{bk:3-14:tani}\R{同}{おな}じデータでも、\R{単位}{たん|い}を\R{変}{か}えたり、\R{全部}{ぜん|ぶ}に\R{同}{おな}じ\R{数}{かず}をたしたりすることがあります。そのとき、\R{平均}{へい|きん}と $\sigma$ がどうなるかを\R{調}{しら}べます。

\bkthrohead{\R{全部}{ぜん|ぶ}に\R{同}{おな}じ\R{数}{かず}をたす}
\R{機械}{き|かい}Aの6\R{本}{ぽん}に、あとから\qty{10}{mL}ずつお\R{茶}{ちゃ}を\R{足}{た}したとします。

\begin{bkcenter}
\begin{bktable}{colspec={Q[l,m,wd=26mm] *{6}{Q[c,m,wd=10.5mm]} Q[c,m,wd=12mm] Q[c,m,wd=10mm]}, rowsep=1pt, colsep=2pt, cells={font=\footnotesize}, row{1}={font=\sffamily\footnotesize, bg=bkTheme!12}}
 & 1 & 2 & 3 & 4 & 5 & 6 & \R{平均}{へい|きん} & $\sigma$ \\
もと（mL） & 501 & 497 & 503 & 500 & 498 & 501 & 500 & 2 \\
\R{偏差}{へん|さ} & $+1$ & $-3$ & $+3$ & 0 & $-2$ & $+1$ & & \\
10\R{足}{た}す（mL） & 511 & 507 & 513 & 510 & 508 & 511 & 510 & 2 \\
\R{偏差}{へん|さ} & $+1$ & $-3$ & $+3$ & 0 & $-2$ & $+1$ & & \\
\end{bktable}
\end{bkcenter}

どの\R{値}{あたい}も10ふえると、\R{平均}{へい|きん}も10ふえて510になります。\R{値}{あたい}も\R{平均}{へい|きん}も10ずつふえたので、\R{偏差}{へん|さ}（\R{値}{あたい}－\R{平均}{へい|きん}）は\R{変}{か}わりません。\R{偏差}{へん|さ}が\R{同}{おな}じなので、$\sigma$ も2のままです。\R{点}{てん}の\R{図}{ず}でいうと、\R{全体}{ぜん|たい}が\R{右}{みぎ}へ10\R{動}{うご}くだけで、ちらばり\R{方}{かた}は\R{同}{おな}じです。

\begin{bkfigure}
\begin{tikzpicture}[x=1mm, y=1mm]
  % 1 mL を 5 mm、495 が 0
  \bkthroaxis{0}{5}{495}{516}{1}{5}{0}
  \foreach \v in {501,497,503,500,498} {\fill[bkBlue] ({(\v-495)*5},2.6) circle (1.1);}
  \fill[bkBlue] ({(501-495)*5},5.2) circle (1.1);
  \foreach \v in {511,507,513,510,508} {\fill[bkOrange] ({(\v-495)*5},2.6) circle (1.1);}
  \fill[bkOrange] ({(511-495)*5},5.2) circle (1.1);
  \draw[bkGray, line width=0.9pt, -{Stealth[length=4pt]}] (25,9) -- (75,9);
  \node[anchor=south, font=\sffamily\footnotesize, inner sep=0pt] at (50,10.2) {\R{全体}{ぜん|たい}が10\R{右}{みぎ}へ};
  \draw[bkBlue!80!black, line width=0.6pt, {Stealth[length=2.4pt]}-{Stealth[length=2.4pt]}] (10,-12) -- (40,-12);
  \draw[bkOrange!80!black, line width=0.6pt, {Stealth[length=2.4pt]}-{Stealth[length=2.4pt]}] (60,-12) -- (90,-12);
  \node[anchor=north, font=\sffamily\footnotesize, inner sep=0pt] at (50,-14) {ちらばりのはばは\R{同}{おな}じ};
\end{tikzpicture}
\bkfigcaption{\R{全部}{ぜん|ぶ}に10をたすと、\R{平均}{へい|きん}は10ふえるが、$\sigma$ は\R{変}{か}わらない}
\end{bkfigure}

\bkthrohead{\R{単位}{たん|い}を\R{変}{か}える：cmからmmへ}
5\R{本}{ほん}のくぎの\R{長}{なが}さを\R{考}{かんが}えます。cmでは4.7、5.1、5.0、5.3、\qty{4.9}{cm}で、\R{平均}{へい|きん}は\qty{5.0}{cm}です。mmになおすと、\qty{10}{mm/cm}をかけて、47、51、50、53、\qty{49}{mm}になります（2-10）。

\begin{bkcenter}
\begin{bktable}{colspec={Q[l,m,wd=30mm] *{5}{Q[c,m,wd=11mm]} Q[c,m,wd=12mm] Q[c,m,wd=12mm] Q[c,m,wd=10mm]}, rowsep=1pt, colsep=2pt, cells={font=\footnotesize}, row{1}={font=\sffamily\footnotesize, bg=bkTheme!12}}
 & \SetCell[c=5]{c} \R{偏差}{へん|さ} & & & & & \R{分散}{ぶん|さん} & $\sigma$ & \\
cmのとき & $-0.3$ & $+0.1$ & 0 & $+0.3$ & $-0.1$ & 0.04 & 0.2 & \unit{cm} \\
mmのとき & $-3$ & $+1$ & 0 & $+3$ & $-1$ & 4 & 2 & \unit{mm} \\
\end{bktable}
\bktabcaption{cmの\R{分散}{ぶん|さん}：$(0.09+0.01+0+0.09+0.01)\div5=0.04$、$\sigma=\sqrt{0.04}=0.2$}
\end{bkcenter}

mmにすると、どの\R{値}{あたい}も10\R{倍}{ばい}、\R{平均}{へい|きん}も10\R{倍}{ばい}（50）なので、\R{偏差}{へん|さ}も10\R{倍}{ばい}になります。\R{偏差}{へん|さ}を2\R{乗}{じょう}すると $10\times10=100$ \R{倍}{ばい}なので、\R{分散}{ぶん|さん}は100\R{倍}{ばい}です（0.04→4）。$\sqrt{100}=10$ なので、$\sigma$ は10\R{倍}{ばい}にもどります（0.2→2）。\R{長}{なが}さそのものは\R{変}{か}わっていないので、\qty{0.2}{cm}と\qty{2}{mm}は\R{同}{おな}じずれです。

\begin{bkcenter}
\begin{bktable}{colspec={Q[l,m,wd=44mm] Q[c,m,wd=34mm] Q[c,m,wd=34mm]}, rowsep=1.2pt, colsep=3pt, cells={font=\footnotesize}, row{1}={font=\sffamily\bfseries\footnotesize, bg=bkTheme!12}}
 & \R{平均}{へい|きん} & $\sigma$ \\
\R{全部}{ぜん|ぶ}に $b$ をたす & $b$ だけふえる & \R{変}{か}わらない \\
\R{全部}{ぜん|ぶ}を $a$ \R{倍}{ばい}する（$a>0$） & $a$ \R{倍}{ばい} & $a$ \R{倍}{ばい}（\R{分散}{ぶん|さん}は $a^2$ \R{倍}{ばい}） \\
\end{bktable}
\bktabcaption{\R{単位}{たん|い}を\R{変}{か}えたときのきまり。\R{覚}{おぼ}えなくても、\R{偏差}{へん|さ}がどうなるかを\R{考}{かんが}えれば\R{作}{つく}れる}
\end{bkcenter}

\bksubsection{5}{\R{何回}{なん|かい}もはかって\R{平均}{へい|きん}すると、ずれは $\dfrac{\sigma}{\sqrt{n}}$}

\leavevmode\label{bk:3-14:rootn}3-13では、\R{何回}{なん|かい}もはかって\R{平均}{へい|きん}すると、ほんとうの\R{値}{あたい}に\R{近}{ちか}づくことを\R{見}{み}ました。では、\R{平均}{へい|きん}のずれは、どのくらい\R{小}{ちい}さくなるのでしょう。いちばんかんたんな\R{場合}{ば|あい}で、\R{全部}{ぜん|ぶ}の\R{場合}{ば|あい}を\R{数}{かぞ}えて\R{調}{しら}べます。

\bkthrohead{いちばんかんたんなはかり\R{方}{かた}}
1\R{回}{かい}はかるたびに、ほんとうの\R{値}{あたい}より「1\R{大}{おお}きい（$+1$）」か「1\R{小}{ちい}さい（$-1$）」かの、どちらかのずれが\R{出}{で}るとします。どちらも\R{同}{おな}じくらい\R{出}{で}ます。1\R{回}{かい}だけのずれは $+1$ か $-1$ なので、2\R{乗}{じょう}はどちらも1、\R{平均}{へい|きん}も1、$\sigma=\sqrt{1}=1$ です。

\noindent\textbf{2\R{回}{かい}はかって\R{平均}{へい|きん}すると}\quad 2\R{回}{かい}のずれの\R{組}{く}み\R{合}{あ}わせは、\R{次}{つぎ}の4\R{通}{とお}りで、どれも\R{同}{おな}じくらい\R{起}{お}こります。

\begin{bkcenter}
\begin{bktable}{colspec={Q[c,m,wd=20mm] Q[c,m,wd=20mm] Q[c,m,wd=34mm] Q[c,m,wd=26mm]}, rowsep=1pt, colsep=3pt, cells={font=\footnotesize}, row{1}={font=\sffamily\footnotesize, bg=bkTheme!12}}
1\R{回}{かい}め & 2\R{回}{かい}め & 2\R{回}{かい}の\R{平均}{へい|きん}のずれ & その2\R{乗}{じょう} \\
$+1$ & $+1$ & $(+1+1)\div2=+1$ & 1 \\
$+1$ & $-1$ & $(+1-1)\div2=0$ & 0 \\
$-1$ & $+1$ & $(-1+1)\div2=0$ & 0 \\
$-1$ & $-1$ & $(-1-1)\div2=-1$ & 1 \\
\end{bktable}
\bktabcaption{\R{平均}{へい|きん}すると、\R{半分}{はん|ぶん}の\R{場合}{ば|あい}でずれが\R{打}{う}ち\R{消}{け}しあって0になる}
\end{bkcenter}

\R{平均}{へい|きん}のずれの2\R{乗}{じょう}をならすと $(1+0+0+1)\div4=\dfrac{2}{4}=\dfrac{1}{2}$。$\sigma$ は $\sqrt{\dfrac{1}{2}}=\dfrac{1}{\sqrt{2}}$ で、\R{約}{やく}0.71です。1\R{回}{かい}のときの1より\R{小}{ちい}さくなりました。

\noindent\textbf{4\R{回}{かい}はかって\R{平均}{へい|きん}すると}\quad 4\R{回}{かい}の\R{組}{く}み\R{合}{あ}わせは $2\times2\times2\times2=16$ \R{通}{とお}りです。$+1$ が\R{何回}{なん|かい}\R{出}{で}たかで\R{分}{わ}けて\R{数}{かぞ}えます。

\begin{bkcenter}
\begin{bktable}{colspec={Q[c,m,wd=16mm] Q[l,m,wd=50mm] Q[c,m,wd=12mm] Q[c,m,wd=20mm] Q[c,m,wd=14mm]}, rowsep=1pt, colsep=2pt, cells={font=\footnotesize}, row{1}={font=\sffamily\footnotesize, bg=bkTheme!12}}
$+1$ の\R{回数}{かい|すう} & \R{並}{なら}び\R{方}{かた}（＋は $+1$、－は $-1$） & \R{通}{とお}り & \R{平均}{へい|きん}のずれ & その2\R{乗}{じょう} \\
4 & ＋＋＋＋ & 1 & $+1$ & 1 \\
3 & －＋＋＋、＋－＋＋、＋＋－＋、＋＋＋－ & 4 & $+0.5$ & 0.25 \\
2 & ＋＋－－、＋－＋－、＋－－＋、－＋＋－、－＋－＋、－－＋＋ & 6 & 0 & 0 \\
1 & ＋－－－、－＋－－、－－＋－、－－－＋ & 4 & $-0.5$ & 0.25 \\
0 & －－－－ & 1 & $-1$ & 1 \\
\end{bktable}
\bktabcaption{16\R{通}{とお}りの\R{並}{なら}び\R{方}{かた}。たとえば $+1$ が3\R{回}{かい}なら、ずれの\R{合計}{ごう|けい}は $3-1=2$、\R{平均}{へい|きん}は $2\div4=0.5$}
\end{bkcenter}

16\R{通}{とお}り\R{全部}{ぜん|ぶ}の2\R{乗}{じょう}をたすと、$1\times1+0.25\times4+0\times6+0.25\times4+1\times1=4$。ならすと $4\div16=\dfrac{1}{4}$ で、$\sigma=\sqrt{\dfrac{1}{4}}=\dfrac{1}{2}$ です。

\begin{bkcenter}
\begin{bktable}{colspec={Q[l,m,wd=40mm] Q[c,m,wd=22mm] Q[c,m,wd=22mm] Q[c,m,wd=22mm]}, rowsep=1.2pt, colsep=3pt, cells={font=\footnotesize}, row{1}={font=\sffamily\footnotesize, bg=bkTheme!12}}
\R{何回}{なん|かい}の\R{平均}{へい|きん}か（$n$） & 1 & 2 & 4 \\
\R{平均}{へい|きん}のずれの $\sigma$ & 1 & $\dfrac{1}{\sqrt{2}}$（\R{約}{やく}0.71） & $\dfrac{1}{2}=\dfrac{1}{\sqrt{4}}$ \\
\end{bktable}
\bktabcaption{$n$ \R{回}{かい}の\R{平均}{へい|きん}のずれは、1\R{回}{かい}のずれを $\sqrt{n}$ でわった\R{大}{おお}きさ}
\end{bkcenter}

2\R{回}{かい}なら $\sqrt{2}$ でわり、4\R{回}{かい}なら $\sqrt{4}=2$ でわった\R{大}{おお}きさになっています。$n$ \R{回}{かい}はかって\R{平均}{へい|きん}すると、
\[ \text{\R{平均}{へい|きん}のずれ（\R{平均}{へい|きん}の $\sigma$）}=\frac{\sigma}{\sqrt{n}} \]
になります。この\R{本}{ほん}では\R{全部}{ぜん|ぶ}の\R{場合}{ば|あい}を\R{数}{かぞ}えて\R{確}{たし}かめましたが、ずれの\R{出}{で}かたがどんな\R{形}{かたち}でも\R{成}{な}り\R{立}{た}つことが\R{知}{し}られています（この\R{表}{ひょう}の1、4、6、4、1は、3-17のパスカルの\R{三角形}{さん|かく|けい}にも\R{出}{で}てきます）。

\bkthrohead{ほんとうのデータで\R{確}{たし}かめる}
3-13のふりこの20\R{回}{かい}を、はかった\R{順}{じゅん}に4\R{回}{かい}ずつ5\R{組}{くみ}に\R{分}{わ}けて、\R{組}{くみ}ごとに\R{平均}{へい|きん}します。

\begin{bkcenter}
\begin{bktable}{colspec={Q[l,m,wd=30mm] *{5}{Q[c,m,wd=15mm]}}, rowsep=1pt, colsep=2pt, cells={font=\footnotesize}, row{1}={font=\sffamily\footnotesize, bg=bkTheme!12}}
\R{組}{くみ} & 1〜4\R{回}{かい}め & 5〜8\R{回}{かい}め & 9〜12\R{回}{かい}め & 13〜16\R{回}{かい}め & 17〜20\R{回}{かい}め \\
4\R{回}{かい}の\R{平均}{へい|きん}（\R{秒}{びょう}） & 14.255 & 14.215 & 14.130 & 14.210 & 14.190 \\
\end{bktable}
\end{bkcenter}

\begin{bkfigure}
\begin{tikzpicture}[x=1mm, y=1mm]
  \bkthroaxis{0}{200}{13.90}{14.50}{0.05}{0.1}{2}
  \foreach \v in {14.37,14.29,14.24,14.12,14.21,14.17,14.33,14.15,14.22,14.19,13.95,14.18,14.23,14.27,14.07,14.20,14.44,14.05} {\fill[bkGray!60] ({(\v-13.9)*200},2.6) circle (1);}
  \fill[bkGray!60] ({(14.16-13.9)*200},2.6) circle (1) ({(14.16-13.9)*200},5.0) circle (1);
  \foreach \v in {14.255,14.215,14.13,14.21,14.19} {\fill[bkBlue] ({(\v-13.9)*200},12) circle (1.3);}
  \node[anchor=west, font=\sffamily\footnotesize, text=bkGray!90!black, inner sep=0pt] at (95,7.6) {1\R{回}{かい}ずつ（$\sigma$ \R{約}{やく}0.11）};
  \node[anchor=west, font=\sffamily\footnotesize, text=bkBlue!80!black, inner sep=0pt] at (68,15.6) {4\R{回}{かい}の\R{平均}{へい|きん}（$\sigma$ \R{約}{やく}0.04）};
\end{tikzpicture}
\bkfigcaption{4\R{回}{かい}ずつ\R{平均}{へい|きん}した5つの\R{点}{てん}（\R{青}{あお}）は、1\R{回}{かい}ずつの\R{点}{てん}（\R{灰色}{はい|いろ}）よりずっとまとまっている}
\end{bkfigure}

5つの\R{平均}{へい|きん}の $\sigma$ を\R{計算}{けい|さん}すると、\R{約}{やく}0.04\R{秒}{びょう}です。1\R{回}{かい}ずつの $\sigma$（\R{約}{やく}0.11\R{秒}{びょう}）よりずっと\R{小}{ちい}さく、$\dfrac{0.11}{\sqrt{4}}=0.055$ と、だいたい\R{合}{あ}っています（\R{組}{くみ}が5つしかないので、ぴったりにはなりません）。

\bkthrohead{3-13のはてな？に、\R{数}{かず}で\R{答}{こた}える}
ふりこを20\R{回}{かい}はかった\R{平均}{へい|きん}14.20\R{秒}{びょう}のずれは、
\[ \frac{\sigma}{\sqrt{20}}=\frac{0.108}{\sqrt{20}}=\frac{0.108}{4.47\cdots}=0.024\cdots \]
で、およそ0.02\R{秒}{びょう}です。だから、20\R{回}{かい}の\R{平均}{へい|きん}は「14.20\R{秒}{びょう}$\pm$0.02\R{秒}{びょう}（\R{平均}{へい|きん}の $\sigma$）」と\R{書}{か}けます。ほんとうの\R{値}{あたい}14.19\R{秒}{びょう}とのちがいは0.01\R{秒}{びょう}で、このはばの\R{中}{なか}に\R{入}{はい}っています。

\begin{bkcenter}
\begin{bktable}{colspec={Q[l,m,wd=36mm] *{6}{Q[c,m,wd=12mm]}}, rowsep=1pt, colsep=2pt, cells={font=\footnotesize}, row{1}={font=\sffamily\footnotesize, bg=bkTheme!12}}
\R{何回}{なん|かい}はかるか（$n$） & 1 & 4 & 9 & 16 & 25 & 100 \\
$\sqrt{n}$ & 1 & 2 & 3 & 4 & 5 & 10 \\
\R{平均}{へい|きん}のずれ & $\sigma$ & $\dfrac{\sigma}{2}$ & $\dfrac{\sigma}{3}$ & $\dfrac{\sigma}{4}$ & $\dfrac{\sigma}{5}$ & $\dfrac{\sigma}{10}$ \\
\end{bktable}
\bktabcaption{ずれを\R{半分}{はん|ぶん}にするには4\R{倍}{ばい}、10\R{分}{ぶん}の1にするには100\R{倍}{ばい}の\R{回数}{かい|すう}がいる}
\end{bkcenter}

ずれを\R{半分}{はん|ぶん}にするのに、\R{回数}{かい|すう}は2\R{倍}{ばい}ではなく4\R{倍}{ばい}もいります。だから、\R{回数}{かい|すう}をふやすだけでなく、3-13の2\R{組}{くみ}のように、はかり\R{方}{かた}をくふうして1\R{回}{かい}の $\sigma$ を\R{小}{ちい}さくすることも\R{大切}{たい|せつ}です。

\begin{seikaku}
\noindent\textbf{$n$ でわるか、$n-1$ でわるか}\par
この\R{章}{しょう}では、\R{分散}{ぶん|さん}を「\R{偏差}{へん|さ}の2\R{乗}{じょう}の\R{合計}{ごう|けい}÷\R{個数}{こ|すう} $n$」で\R{求}{もと}めました。ところが、\R{大学}{だい|がく}の\R{本}{ほん}や\R{表計算}{ひょう|けい|さん}では、$n$ ではなく $n-1$ でわることもあります（\R{表計算}{ひょう|けい|さん}の \texttt{STDEV.P} は $n$ でわり、\texttt{STDEV.S} は $n-1$ でわります）。

\R{偏差}{へん|さ}は、データから\R{計算}{けい|さん}した\R{平均}{へい|きん}を\R{使}{つか}います。この\R{平均}{へい|きん}は、そのデータにいちばんぴったり\R{合}{あ}うように\R{決}{き}まっているので、ほんとうの\R{値}{あたい}を\R{使}{つか}ったときより、\R{偏差}{へん|さ}の2\R{乗}{じょう}の\R{合計}{ごう|けい}が\R{少}{すこ}し\R{小}{ちい}さく\R{出}{で}ます。その\R{分}{ぶん}をおぎなうために、1つ\R{少}{すく}ない $n-1$ でわるのです。\R{機械}{き|かい}Aなら $24\div5=4.8$、$\sqrt{4.8}=2.19\cdots$ です。データが\R{多}{おお}ければ、$n$ と $n-1$ の\R{差}{さ}はほとんどなくなります。この\R{本}{ほん}では、$n$ でわる\R{求}{もと}め\R{方}{かた}を\R{使}{つか}います。
\end{seikaku}

\begin{taisetsu}
\R{偏差}{へん|さ}＝\R{値}{あたい}－\R{平均}{へい|きん}。\R{偏差}{へん|さ}を\R{全部}{ぜん|ぶ}たすと、かならず0。\\
\R{分散}{ぶん|さん}＝\R{偏差}{へん|さ}の2\R{乗}{じょう}の\R{平均}{へい|きん}。\R{標準偏差}{ひょう|じゅん|へん|さ} $\sigma=\sqrt{\text{\R{分散}{ぶん|さん}}}$（もとの\R{単位}{たん|い}にもどる）。\\
つりがね\R{形}{がた}なら、\R{平均}{へい|きん}$\pm\sigma$ にだいたい7\R{割}{わり}、$\pm2\sigma$ にだいたい95\%。\\
\R{全部}{ぜん|ぶ}に $b$ をたしても $\sigma$ は\R{同}{おな}じ。$a$ \R{倍}{ばい}すると $\sigma$ も $a$ \R{倍}{ばい}。\\
$n$ \R{回}{かい}の\R{平均}{へい|きん}のずれは $\dfrac{\sigma}{\sqrt{n}}$。
\end{taisetsu}

\begin{kotoba}
\small
\bkword{\R{偏差}{へん|さ}}{\R{値}{あたい}－\R{平均}{へい|きん}。\R{平均}{へい|きん}からのずれ。たすと0。}
\bkword{\R{分散}{ぶん|さん}}{\R{偏差}{へん|さ}の2\R{乗}{じょう}の\R{平均}{へい|きん}。\R{単位}{たん|い}はもとの\R{単位}{たん|い}の2\R{乗}{じょう}。}
\bkword{\R{標準偏差}{ひょう|じゅん|へん|さ} $\sigma$}{\R{分散}{ぶん|さん}の $\sqrt{\phantom{2}}$。「だいたいこのくらいずれる」を\R{表}{あらわ}す\R{数}{かず}。}
\bkword{\R{平均}{へい|きん}$\pm\sigma$}{\R{平均}{へい|きん}より $\sigma$ \R{小}{ちい}さい\R{値}{あたい}から、$\sigma$ \R{大}{おお}きい\R{値}{あたい}までの\R{範囲}{はん|い}。}
\end{kotoba}

%% ---------------- 例題 ----------------
\bksection{reidai}{\R{例題}{れい|だい}でたしかめよう}

\begin{reidai}[\R{偏差}{へん|さ}を\R{求}{もと}める]
ばねに\R{同}{おな}じおもりをつるして、のびを5\R{回}{かい}はかると、42、45、40、44、\qty{39}{mm}でした。それぞれの\R{偏差}{へん|さ}を\R{求}{もと}め、\R{偏差}{へん|さ}の\R{合計}{ごう|けい}が0になることを\R{確}{たし}かめましょう。\par
\kangaekata まず\R{平均}{へい|きん}を\R{求}{もと}めます。\R{偏差}{へん|さ}は「\R{値}{あたい}－\R{平均}{へい|きん}」です。\R{平均}{へい|きん}より\R{小}{ちい}さい\R{値}{あたい}の\R{偏差}{へん|さ}はマイナスです。\par
\shiki ① \R{平均}{へい|きん}：$42+45+40+44+39=210$、\ $210\div5=42$\par
\hspace*{2.6\zw}② \R{偏差}{へん|さ}：$42-42=0$、\ $45-42=+3$、\ $40-42=-2$\par
\hspace*{2.6\zw}\phantom{②}\ $44-42=+2$、\ $39-42=-3$\par
\hspace*{2.6\zw}③ \R{合計}{ごう|けい}：$0+3+(-2)+2+(-3)=0$\par
\kotae \R{偏差}{へん|さ}は0、$+3$、$-2$、$+2$、$-3$（mm）。\R{合計}{ごう|けい}は0
\end{reidai}

\begin{jibun}
ペットボトルロケットを5\R{回}{かい}とばすと、63、58、61、66、\qty{57}{m}でした。それぞれの\R{偏差}{へん|さ}を\R{求}{もと}め、\R{合計}{ごう|けい}が0になることを\R{確}{たし}かめましょう。
\bkjans{\R{平均}{へい|きん} $305\div5=61$。\R{偏差}{へん|さ}は $+2$、$-3$、0、$+5$、$-4$（m）。\R{合計}{ごう|けい} $2-3+0+5-4=0$}
\end{jibun}

\begin{reidai}[\R{分散}{ぶん|さん}と\R{標準偏差}{ひょう|じゅん|へん|さ}]
5\R{本}{ほん}のくぎの\R{長}{なが}さをはかると、51、47、50、53、\qty{49}{mm}でした。\R{分散}{ぶん|さん}と\R{標準偏差}{ひょう|じゅん|へん|さ} $\sigma$ を\R{求}{もと}めましょう。\par
\kangaekata ①\R{平均}{へい|きん}→②\R{偏差}{へん|さ}→③2\R{乗}{じょう}→④\R{平均}{へい|きん}（\R{分散}{ぶん|さん}）→⑤ $\sqrt{\phantom{2}}$（$\sigma$）の\R{順}{じゅん}に\R{計算}{けい|さん}します。\R{表}{ひょう}にすると\R{見}{み}やすくなります。\par
\shiki ① \R{平均}{へい|きん}：$(51+47+50+53+49)\div5=250\div5=50$\par
\begin{center}
\begin{bktable}{colspec={Q[l,m,wd=26mm] *{5}{Q[c,m,wd=11mm]} Q[c,m,wd=12mm]}, rowsep=0.6pt, colsep=2pt, cells={font=\footnotesize}, row{1}={font=\sffamily\footnotesize, bg=white}}
\R{長}{なが}さ（mm） & 51 & 47 & 50 & 53 & 49 & \R{合計}{ごう|けい} \\
② \R{偏差}{へん|さ} & $+1$ & $-3$ & 0 & $+3$ & $-1$ & 0 \\
③ \R{偏差}{へん|さ}の2\R{乗}{じょう} & 1 & 9 & 0 & 9 & 1 & 20 \\
\end{bktable}
\end{center}
\hspace*{2.6\zw}④ \R{分散}{ぶん|さん}：$20\div5=4$\par
\hspace*{2.6\zw}⑤ $\sigma=\sqrt{4}=2$\par
\kotae \R{分散}{ぶん|さん}4、\R{標準偏差}{ひょう|じゅん|へん|さ}\qty{2}{mm}
\end{reidai}

\begin{jibun}
5\R{回}{かい}の\R{記録}{き|ろく}が、26、14、20、22、18\R{点}{てん}でした。\R{分散}{ぶん|さん}と\R{標準偏差}{ひょう|じゅん|へん|さ} $\sigma$ を\R{求}{もと}めましょう。
\bkjans{\R{平均}{へい|きん}20。\R{偏差}{へん|さ} $+6$、$-6$、0、$+2$、$-2$。2\R{乗}{じょう}36、36、0、4、4の\R{合計}{ごう|けい}80。\R{分散}{ぶん|さん} $80\div5=16$、$\sigma=\sqrt{16}=4$ で4\R{点}{てん}}
\end{jibun}

\begin{reidai}[$\sigma$ でばらつきをくらべる]
2\R{台}{だい}のロボットカーXとYを5\R{回}{かい}ずつ\R{走}{はし}らせ、\R{止}{と}まった\R{所}{ところ}までの\R{距離}{きょ|り}をはかりました。どちらのほうが、\R{止}{と}まる\R{所}{ところ}がそろっていますか。$\sigma$ をくらべましょう。\par
\hspace*{1\zw}X：6、8、7、9、\qty{5}{m}\qquad Y：3、10、7、11、\qty{4}{m}\par
\kangaekata どちらも\R{平均}{へい|きん}は\qty{7}{m}です。$\sqrt{\phantom{2}}$ が\R{割}{わ}り\R{切}{き}れないときは、\R{電卓}{でん|たく}の $\sqrt{\phantom{2}}$ キーを\R{使}{つか}います。\par
\shiki X\enspace \R{偏差}{へん|さ} $-1$、$+1$、0、$+2$、$-2$。2\R{乗}{じょう}の\R{合計}{ごう|けい} $1+1+0+4+4=10$\par
\hspace*{2.6\zw}\phantom{X}\enspace \R{分散}{ぶん|さん} $10\div5=2$、\ $\sigma=\sqrt{2}=1.41\cdots$\par
\hspace*{2.6\zw}Y\enspace \R{偏差}{へん|さ} $-4$、$+3$、0、$+4$、$-3$。2\R{乗}{じょう}の\R{合計}{ごう|けい} $16+9+0+16+9=50$\par
\hspace*{2.6\zw}\phantom{Y}\enspace \R{分散}{ぶん|さん} $50\div5=10$、\ $\sigma=\sqrt{10}=3.16\cdots$\par
\kotae Xのほうがそろっている（$\sigma$ はX\R{約}{やく}\qty{1.4}{m}、Y\R{約}{やく}\qty{3.2}{m}）
\end{reidai}

\begin{jibun}
2\R{人}{にん}が\R{同}{おな}じ\R{的}{まと}に5\R{回}{かい}ずつボールを\R{投}{な}げ、\R{的}{まと}の\R{点数}{てん|すう}を\R{記録}{き|ろく}しました。$\sigma$ をくらべて、\R{点数}{てん|すう}がそろっているほうを\R{答}{こた}えましょう。\par
P：21、19、20、23、17\R{点}{てん}\qquad Q：25、15、20、18、22\R{点}{てん}
\bkjans{どちらも\R{平均}{へい|きん}20。P：2\R{乗}{じょう}の\R{合計}{ごう|けい}20、\R{分散}{ぶん|さん}4、$\sigma=2$。Q：2\R{乗}{じょう}の\R{合計}{ごう|けい}58、\R{分散}{ぶん|さん}11.6、$\sigma=\sqrt{11.6}=3.40\cdots$。Pのほうがそろっている}
\end{jibun}

\begin{reidai}[\R{平均}{へい|きん}$\pm\sigma$ に\R{入}{はい}る\R{数}{かず}]
ある\R{乾電池}{かん|でん|ち}でおもちゃを\R{動}{うご}かし、\R{止}{と}まるまでの\R{時間}{じ|かん}を10\R{回}{かい}はかりました。\R{平均}{へい|きん}は50\R{分}{ぷん}、$\sigma$ は3\R{分}{ぷん}でした。（1）\R{平均}{へい|きん}$\pm\sigma$ の\R{範囲}{はん|い}を\R{求}{もと}めましょう。（2）その\R{範囲}{はん|い}に\R{入}{はい}るのは\R{何回}{なん|かい}で、\R{全体}{ぜん|たい}の\R{何}{なん}\%ですか。\par
\hspace*{1\zw}51、46、53、49、54、50、44、51、49、53（\R{分}{ふん}）\par
\kangaekata 「\R{平均}{へい|きん}$\pm\sigma$」は、\R{平均}{へい|きん}$-\sigma$ から\R{平均}{へい|きん}$+\sigma$ までです。\R{範囲}{はん|い}のはしの\R{値}{あたい}も\R{中}{なか}に\R{入}{い}れて\R{数}{かぞ}えます。\par
\shiki（1）$50-3=47$、\ $50+3=53$ で、47\R{分}{ふん}から53\R{分}{ぷん}まで\par
\hspace*{2.6\zw}（2）\R{入}{はい}るのは 51、53、49、50、51、49、53 の7\R{回}{かい}\par
\hspace*{2.6\zw}\phantom{（2）}\R{入}{はい}らないのは 46、54、44 の3\R{回}{かい}。$7\div10=0.7$\par
\begin{center}
\begin{tikzpicture}[x=1mm, y=1mm]
  % 1分を7mm、43分が0
  \fill[bkGreen!25] (28,-0.6) rectangle (70,9.5);
  \bkthroaxis{0}{7}{43}{55}{1}{1}{0}
  \foreach \v/\h in {44/1,46/1,49/1,49/2,50/1,51/1,51/2,53/1,53/2,54/1} {
    \pgfmathsetmacro\c{(\v>=47 && \v<=53) ? 1 : 0}
    \ifdim\c pt>0.5pt \fill[bkBlue] ({(\v-43)*7},{\h*2.6}) circle (1.1); \else \fill[bkRed] ({(\v-43)*7},{\h*2.6}) circle (1.1); \fi}
  \node[anchor=south, font=\sffamily\footnotesize, text=bkGreen!30!black, inner sep=0pt] at (49,10.2) {\R{平均}{へい|きん}$\pm\sigma$（47〜53）：7\R{回}{かい}};
  \node[anchor=west, font=\sffamily\footnotesize, inner sep=0pt] at (88,0) {（\R{分}{ふん}）};
\end{tikzpicture}
\end{center}
\kotae（1）47\R{分}{ふん}から53\R{分}{ぷん}まで\quad（2）7\R{回}{かい}、70\%
\end{reidai}

\begin{jibun}
10\R{回}{かい}はかった\R{値}{あたい}の\R{平均}{へい|きん}は20、$\sigma$ は2でした。（1）\R{平均}{へい|きん}$\pm\sigma$ の\R{範囲}{はん|い}（2）その\R{範囲}{はん|い}に\R{入}{はい}る\R{回数}{かい|すう}と\%を\R{求}{もと}めましょう。\par
19、22、17、21、23、19、17、22、21、19
\bkjans{（1）\nobreak 18から22まで（2）\nobreak 19、22、21、19、22、21、19の7\R{回}{かい}。$7\div10=0.7$ で70\%}
\end{jibun}

\begin{reidai}[\R{単位}{たん|い}を\R{変}{か}えたときの\R{平均}{へい|きん}と $\sigma$]
5\R{本}{ほん}のえんぴつの\R{長}{なが}さは15.1、14.7、15.0、15.3、\qty{14.9}{cm}で、\R{平均}{へい|きん}は\qty{15.0}{cm}、$\sigma$ は\qty{0.2}{cm}です。（1）mmで\R{表}{あらわ}すと、\R{平均}{へい|きん}と $\sigma$ はどうなりますか。（2）どのえんぴつにも、\R{長}{なが}さ\qty{0.5}{cm}のキャップをつけた\R{長}{なが}さでは、どうなりますか。（3）mmで\R{表}{あらわ}してから\qty{5}{mm}のキャップをつけると、どうなりますか。\par
\kangaekata \R{全部}{ぜん|ぶ}を10\R{倍}{ばい}すると、\R{偏差}{へん|さ}も10\R{倍}{ばい}なので $\sigma$ も10\R{倍}{ばい}。\R{全部}{ぜん|ぶ}に\R{同}{おな}じ\R{数}{かず}をたすと、\R{偏差}{へん|さ}は\R{変}{か}わらないので $\sigma$ も\R{変}{か}わりません。\par
\shiki（1）\qty{10}{mm/cm}をかける。\R{平均}{へい|きん} $15.0\times10=150$、\ $\sigma$ $0.2\times10=2$\par
\hspace*{2.6\zw}（2）\R{平均}{へい|きん} $15.0+0.5=15.5$、\ $\sigma$ は0.2のまま\par
\hspace*{2.6\zw}（3）（1）に5をたす。\R{平均}{へい|きん} $150+5=155$、\ $\sigma$ は2のまま\par
\bkstep[bkGreen]{\R{確}{たし}かめ}mmの\R{偏差}{へん|さ}は $+1$、$-3$、0、$+3$、$-1$。\R{分散}{ぶん|さん}は $20\div5=4$、$\sigma=2$ で\R{合}{あ}う。\par
\kotae（1）\R{平均}{へい|きん}\qty{150}{mm}、$\sigma$ \qty{2}{mm}（2）\R{平均}{へい|きん}\qty{15.5}{cm}、$\sigma$ \qty{0.2}{cm}（3）\R{平均}{へい|きん}\qty{155}{mm}、$\sigma$ \qty{2}{mm}
\end{reidai}

\begin{jibun}
\R{機械}{き|かい}Aのお\R{茶}{ちゃ}は、\R{平均}{へい|きん}\qty{500}{mL}、$\sigma$ \qty{2}{mL}でした。（1）Lで\R{表}{あらわ}すと、\R{平均}{へい|きん}と $\sigma$ は？（2）どのボトルにも\qty{10}{mL}ずつ\R{足}{た}すと？（3）Lで\R{表}{あらわ}してから\qty{0.01}{L}ずつ\R{足}{た}すと？
\bkjans{（1）\nobreak \qty{1000}{mL/L}でわる。\R{平均}{へい|きん}\qty{0.5}{L}、$\sigma$ \qty{0.002}{L}（2）\nobreak \R{平均}{へい|きん}\qty{510}{mL}、$\sigma$ \qty{2}{mL}（3）\nobreak \R{平均}{へい|きん}\qty{0.51}{L}、$\sigma$ \qty{0.002}{L}}
\end{jibun}

\begin{reidai}[\R{何回}{なん|かい}もはかった\R{平均}{へい|きん}のずれ]
ストップウォッチで1\R{回}{かい}はかるときの $\sigma$ は0.12\R{秒}{びょう}です。（1）4\R{回}{かい}はかって\R{平均}{へい|きん}すると、\R{平均}{へい|きん}のずれは\R{何秒}{なん|びょう}ですか。（2）36\R{回}{かい}なら\R{何秒}{なん|びょう}ですか。（3）\R{平均}{へい|きん}のずれを0.03\R{秒}{びょう}にするには、\R{何回}{なん|かい}はかればよいですか。\par
\kangaekata ずれは $\sigma\div\sqrt{n}$。（3）は、$\sqrt{n}$ が\R{何}{なん}になればよいかを\R{先}{さき}に\R{考}{かんが}えます。\par
\shiki（1）$\sqrt{4}=2$ なので、$0.12\div2=0.06$\par
\hspace*{2.6\zw}（2）$\sqrt{36}=6$ なので、$0.12\div6=0.02$\par
\hspace*{2.6\zw}（3）0.12を0.03にするには、$0.12\div0.03=4$ でわればよい\par
\hspace*{2.6\zw}\phantom{（3）}$\sqrt{n}=4$ になる $n$ は、$4^2=16$\par
\begin{center}
\begin{tikzpicture}[x=1mm, y=1mm]
  % 0.01秒を6mm
  \foreach \y/\n/\v/\c in {14/{1\R{回}{かい}}/0.12/bkGray, 7/{4\R{回}{かい}}/0.06/bkBlue, 0/{36\R{回}{かい}}/0.02/bkGreen} {
    \fill[\c!45] (0,\y) rectangle ({\v*600},\y+4.6);
    \draw[\c!80!black, line width=0.6pt] (0,\y) rectangle ({\v*600},\y+4.6);
    \node[anchor=east, font=\sffamily\footnotesize, inner sep=0pt] at (-2.4,\y+2.3) {\n};
    \node[anchor=west, font=\sffamily\footnotesize, inner sep=0pt] at ({\v*600+2},\y+2.3) {\v\R{秒}{びょう}};}
  \draw[bkInk, line width=0.7pt] (0,-1.4) -- (0,20);
  \node[anchor=west, font=\sffamily\footnotesize, align=left, inner sep=0pt] at (56,2.6) {\R{回数}{かい|すう}を9\R{倍}{ばい}にすると、\\ずれは3\R{分}{ぶん}の1};
\end{tikzpicture}
\end{center}
\kotae（1）0.06\R{秒}{びょう}\quad（2）0.02\R{秒}{びょう}\quad（3）16\R{回}{かい}
\end{reidai}

\begin{jibun}
1\R{回}{かい}はかるときの $\sigma$ が\qty{0.5}{g}のはかりがあります。（1）25\R{回}{かい}の\R{平均}{へい|きん}のずれは？（2）100\R{回}{かい}なら？（3）\R{平均}{へい|きん}のずれを\qty{0.25}{g}にするには\R{何回}{なん|かい}？
\bkjans{（1）\nobreak $0.5\div\sqrt{25}=0.5\div5=0.1$ で\qty{0.1}{g}（2）\nobreak $0.5\div10=0.05$ で\qty{0.05}{g}（3）\nobreak $0.5\div0.25=2$ なので $\sqrt{n}=2$、$n=4$ で4\R{回}{かい}}
\end{jibun}

\begin{column}{テストの「\R{偏差値}{へん|さ|ち}」も $\sigma$ で\R{決}{き}まる}
\R{模擬試験}{も|ぎ|し|けん}などで\R{使}{つか}われる「\R{偏差値}{へん|さ|ち}」は、\R{自分}{じ|ぶん}の\R{点}{てん}の\R{偏差}{へん|さ}が $\sigma$ の\R{何倍}{なん|ばい}かを\R{使}{つか}って、\R{次}{つぎ}のように\R{決}{き}めた\R{数}{かず}です。
\begin{center}
\R{偏差値}{へん|さ|ち}＝50＋10×（\R{自分}{じ|ぶん}の\R{点}{てん}－\R{平均}{へい|きん}）÷ $\sigma$
\end{center}
\R{平均}{へい|きん}60\R{点}{てん}、$\sigma$ が10\R{点}{てん}のテストで80\R{点}{てん}なら、\R{偏差}{へん|さ}は20で $\sigma$ の2\R{倍}{ばい}なので、\R{偏差値}{へん|さ|ち}は $50+10\times2=70$ です。\R{平均}{へい|きん}ちょうどなら50。\R{点}{てん}のちらばりが\R{大}{おお}きいテストでも\R{小}{ちい}さいテストでも、「みんなの\R{中}{なか}でどのあたりにいるか」を\R{同}{おな}じものさしでくらべられるように、$\sigma$ でわっているのです。
\end{column}

%% ---------------- まちがえやすいところ ----------------
\bksection{tsumazuki}{まちがえやすいところをチェック}

\noindent\begin{minipage}[t]{0.485\linewidth}\vspace{0pt}\raggedright
\begin{tsumazuki}{\R{偏差}{へん|さ}の\R{合計}{ごう|けい}で、ばらつきを\R{見}{み}る}
\R{偏差}{へん|さ}の\R{合計}{ごう|けい}が0になったので……
\bkng{ばらつきはない}
\bkok{\R{偏差}{へん|さ}の\R{合計}{ごう|けい}は、いつも0}
プラスとマイナスが\R{打}{う}ち\R{消}{け}しあうだけです。2\R{乗}{じょう}して\R{平均}{へい|きん}し、$\sigma$ でくらべます。
\end{tsumazuki}
\end{minipage}\hfill
\begin{minipage}[t]{0.485\linewidth}\vspace{0pt}\raggedright
\begin{tsumazuki}{マイナスの\R{偏差}{へん|さ}を2\R{乗}{じょう}する}
\R{偏差}{へん|さ}が $-3$ のとき、2\R{乗}{じょう}は……
\bkng{$(-3)^2=-9$}
\bkok{$(-3)^2=(-3)\times(-3)=9$}
マイナス×マイナスはプラスです（2-2）。2\R{乗}{じょう}がマイナスになることはありません。
\end{tsumazuki}
\end{minipage}

\noindent\begin{minipage}[t]{0.485\linewidth}\vspace{0pt}\raggedright
\begin{tsumazuki}{\R{分散}{ぶん|さん}を $\sigma$ と\R{思}{おも}う}
\R{偏差}{へん|さ}の2\R{乗}{じょう}の\R{平均}{へい|きん}が4（mLの2\R{乗}{じょう}）のとき……
\bkng{$\sigma=4$\,mL}
\bkok{$\sigma=\sqrt{4}=2$\,mL}
\R{分散}{ぶん|さん}は\R{単位}{たん|い}も2\R{乗}{じょう}になっています。$\sqrt{\phantom{2}}$ をとって、もとの\R{単位}{たん|い}にもどします。
\end{tsumazuki}
\end{minipage}\hfill
\begin{minipage}[t]{0.485\linewidth}\vspace{0pt}\raggedright
\begin{tsumazuki}{\R{同}{おな}じ\R{数}{かず}をたすと $\sigma$ もふえる}
\R{平均}{へい|きん}500、$\sigma$ 2の\R{全部}{ぜん|ぶ}に10をたすと……
\bkng{$\sigma$ は $2+10=12$}
\bkok{$\sigma$ は2のまま}
\R{値}{あたい}も\R{平均}{へい|きん}も10ふえるので、\R{偏差}{へん|さ}は\R{変}{か}わりません。ちらばり\R{方}{かた}は\R{同}{おな}じです。
\end{tsumazuki}
\end{minipage}

\noindent\begin{minipage}[t]{0.485\linewidth}\vspace{0pt}\raggedright
\begin{tsumazuki}{10\R{倍}{ばい}すると $\sigma$ は100\R{倍}{ばい}}
$\sigma$ \qty{0.2}{cm}をmmにすると……
\bkng{$0.2\times100=20$\,mm}
\bkok{$0.2\times10=2$\,mm}
100\R{倍}{ばい}になるのは\R{分散}{ぶん|さん}です。$\sigma$ は $\sqrt{100}=10$ \R{倍}{ばい}です。
\end{tsumazuki}
\end{minipage}\hfill
\begin{minipage}[t]{0.485\linewidth}\vspace{0pt}\raggedright
\begin{tsumazuki}{\R{平均}{へい|きん}のずれを $n$ でわる}
$\sigma$ 0.12\R{秒}{びょう}で、4\R{回}{かい}の\R{平均}{へい|きん}のずれは……
\bkng{$0.12\div4=0.03$}
\bkok{$0.12\div\sqrt{4}=0.06$}
ずれは\R{全部}{ぜん|ぶ}は\R{打}{う}ち\R{消}{け}しあいません。$n$ ではなく $\sqrt{n}$ でわります。
\end{tsumazuki}
\end{minipage}

%% ---------------- 物理・工学ではこう使う ----------------
\begin{butsuri}
\noindent\begin{minipage}[c]{0.56\linewidth}
\begin{itemize}[leftmargin=1.2\zw, itemsep=1pt, topsep=0pt, label={\color{bkPurple}●}]
\item \textbf{\R{計測}{けい|そく}の\R{不確}{ふ|たし}かさ}\par\noindent \R{理科}{り|か}や\R{工学}{こう|がく}では、はかった\R{結果}{けっ|か}を「14.20\R{秒}{びょう}$\pm$0.02\R{秒}{びょう}」のように、ずれの\R{目安}{め|やす}（\textbf{\R{不確}{ふ|たし}かさ}）といっしょに\R{書}{か}きます。ふつうは\R{平均}{へい|きん}の $\sigma$（$\sigma\div\sqrt{n}$）を\R{使}{つか}います。これがあると、ほかの\R{結果}{けっ|か}とくらべて「ちがう」のか「ずれの\R{中}{なか}で\R{同}{おな}じ」なのかを\R{決}{き}められます。
\end{itemize}
\end{minipage}\hfill
\begin{minipage}[c]{0.41\linewidth}\centering
\begin{tikzpicture}[x=1mm, y=1mm]
  \draw[bkInk, line width=0.6pt] (0,0) -- (50,0);
  \foreach \x/\l in {5/14.15,25/14.20,45/14.25} {\draw[bkInk, line width=0.6pt] (\x,-1) -- (\x,1); \node[anchor=north, font=\sffamily\scriptsize, inner sep=0pt] at (\x,-2) {\l};}
  \fill[bkYellow!25] (17,8.5) rectangle (29,27.5);
  \draw[bkPurple, line width=1pt] (17,11) -- (33,11) (17,9) -- (17,13) (33,9) -- (33,13);
  \fill[bkPurple] (25,11) circle (1.3);
  \node[anchor=south, font=\sffamily\scriptsize, text=bkPurple!80!black, align=center, inner sep=0pt] at (25,14.2) {ソラさんのクラス\\（14.20\R{秒}{びょう}$\pm$0.02\R{秒}{びょう}）};
  \draw[bkGreen!60!black, line width=1pt] (13,25) -- (29,25) (13,23) -- (13,27) (29,23) -- (29,27);
  \fill[bkGreen!60!black] (21,25) circle (1.3);
  \node[anchor=south, font=\sffamily\scriptsize, text=bkGreen!40!black, align=center, inner sep=0pt] at (21,28.6) {となりのクラス\\（14.19\R{秒}{びょう}$\pm$0.02\R{秒}{びょう}）};
  \node[anchor=north, font=\sffamily\scriptsize, align=center, inner sep=0pt] at (25,-6.6) {\R{黄色}{き|いろ}の\R{所}{ところ}で2つの\R{区間}{く|かん}が\R{重}{かさ}なる\\→ずれの\R{中}{なか}で\R{同}{おな}じと\R{考}{かんが}えてよい};
\end{tikzpicture}
\end{minipage}

\par\vspace{4pt}
\noindent\begin{minipage}[c]{0.56\linewidth}
\begin{itemize}[leftmargin=1.2\zw, itemsep=1pt, topsep=0pt, label={\color{bkPurple}●}]
\item \textbf{\R{工場}{こう|じょう}の\R{品質管理}{ひん|しつ|かん|り}}\par\noindent \R{製品}{せい|ひん}をはかった\R{値}{あたい}を\R{順}{じゅん}に\R{点}{てん}でかき、\R{平均}{へい|きん}と\R{平均}{へい|きん}$\pm3\sigma$ の\R{線}{せん}を\R{引}{ひ}いた\textbf{\R{管理図}{かん|り|ず}}を\R{使}{つか}います。つりがね\R{形}{がた}なら、\R{点}{てん}が $\pm3\sigma$ の\R{外}{そと}に\R{出}{で}るのは1000\R{回}{かい}に3\R{回}{かい}ほどです。\R{外}{そと}に\R{出}{で}たら、\R{機械}{き|かい}に\R{何}{なに}か\R{起}{お}きたと\R{考}{かんが}えて\R{調}{しら}べます。
\end{itemize}
\end{minipage}\hfill
\begin{minipage}[c]{0.41\linewidth}\centering
\begin{tikzpicture}[x=0.82mm, y=1mm]
  % 1 mL を 2 mm、492 が 0。番号1を2.4mm
  \draw[bkInk, line width=0.6pt] (0,0) -- (52,0) (0,0) -- (0,34);
  \draw[bkGreen!60!black, line width=0.8pt] (0,16) -- (51,16);
  \draw[bkRed, line width=0.8pt, dash pattern=on 2pt off 1pt] (0,4) -- (51,4) (0,28) -- (51,28);
  \foreach \k/\v in {1/501,2/499,3/500,4/502,5/498,6/500,7/503,8/497,9/501,10/499,11/500,12/502,13/498,14/501,15/500,16/499,17/507,18/501,19/500,20/499} {\fill[bkPurple] (\k*2.4,{(\v-492)*2}) circle (0.8mm);}
  \draw[bkPurple, line width=0.4pt] plot coordinates {(2.4,18) (4.8,14) (7.2,16) (9.6,20) (12,12) (14.4,16) (16.8,22) (19.2,10) (21.6,18) (24,14) (26.4,16) (28.8,20) (31.2,12) (33.6,18) (36,16) (38.4,14) (40.8,30) (43.2,18) (45.6,16) (48,14)};
  \draw[bkRed, line width=0.7pt] (40.8,30) circle (2mm);
  \node[anchor=west, font=\sffamily\scriptsize, text=bkRed!80!black, inner sep=0pt] at (52,28) {$+3\sigma$};
  \node[anchor=west, font=\sffamily\scriptsize, text=bkGreen!50!black, inner sep=0pt] at (52,16) {\R{平均}{へい|きん}};
  \node[anchor=west, font=\sffamily\scriptsize, text=bkRed!80!black, inner sep=0pt] at (52,4) {$-3\sigma$};
  \node[anchor=north, font=\sffamily\scriptsize, inner sep=0pt] at (26,-1.6) {\R{作}{つく}った\R{順}{じゅん}};
\end{tikzpicture}
\end{minipage}

\par\vspace{4pt}
\noindent\begin{minipage}[c]{0.72\linewidth}
\begin{itemize}[leftmargin=1.2\zw, itemsep=1pt, topsep=0pt, label={\color{bkPurple}●}]
\item \textbf{\R{放射線}{ほう|しゃ|せん}を\R{数}{かぞ}える}\par\noindent ガイガーカウンターは、\R{放射線}{ほう|しゃ|せん}が1つ\R{通}{とお}るたびに1つ\R{数}{かぞ}えます。たまにしか\R{起}{お}きないできごとを\R{数}{かぞ}えると、\R{数}{かぞ}えた\R{数}{かず} $N$ の $\sigma$ はおよそ $\sqrt{N}$ です。1\R{分}{ぷん}で100なら $\pm10$（10\%）。10\R{分}{ぷん}で1000なら\R{約}{やく} $\pm32$ で、1\R{分}{ぷん}あたり $100\pm3.2$（3.2\%）。\R{長}{なが}く\R{数}{かぞ}えるほどずれの\R{割合}{わり|あい}が\R{小}{ちい}さくなるのは、⑤と\R{同}{おな}じきまりです。
\end{itemize}
\end{minipage}\hfill
\begin{minipage}[c]{0.25\linewidth}\centering
\begin{tikzpicture}[x=1mm, y=1mm]
  \filldraw[fill=bkYellow!30, draw=bkInk!70, line width=0.7pt, rounded corners=2mm] (0,0) rectangle (28,20);
  \bkdisplay{(2,7)}{24}{100}
  \node[anchor=north, font=\sffamily\scriptsize, inner sep=0pt] at (14,5.6) {1\R{分間}{ふん|かん}に\R{数}{かぞ}えた\R{数}{かず}};
\end{tikzpicture}
\end{minipage}
\end{butsuri}

%% ---------------- 章末問題 ----------------
\bkendproblems
\bkproblemgroup{\R{基本}{き|ほん}}{3}

\bkproblem 50m\R{走}{そう}を5\R{回}{かい}はかると、8.4、8.1、8.3、8.0、8.2\R{秒}{びょう}でした。それぞれの\R{偏差}{へん|さ}を\R{求}{もと}め、\R{合計}{ごう|けい}が0になることを\R{確}{たし}かめましょう。
\bkpans{\R{平均}{へい|きん} $41.0\div5=8.2$。\R{偏差}{へん|さ}は $+0.2$、$-0.1$、$+0.1$、$-0.2$、0（\R{秒}{びょう}）。\R{合計}{ごう|けい}0}

\bkproblem 5\R{回}{かい}の\R{記録}{き|ろく}が33、27、30、31、29でした。\R{分散}{ぶん|さん}と\R{標準偏差}{ひょう|じゅん|へん|さ} $\sigma$ を\R{求}{もと}めましょう。
\bkpans{\R{平均}{へい|きん}30。\R{偏差}{へん|さ} $+3$、$-3$、0、$+1$、$-1$。2\R{乗}{じょう}の\R{合計}{ごう|けい}20。\R{分散}{ぶん|さん}4、$\sigma=2$}

\bkproblem 2つの\R{温度計}{おん|ど|けい}SとTで、\R{同}{おな}じ\R{部屋}{へ|や}の\R{温度}{おん|ど}を5\R{回}{かい}ずつはかりました。$\sigma$ をくらべて、\R{値}{あたい}がそろっているほうを\R{答}{こた}えましょう。\par
S：21、23、22、24、\qty{20}{\degreeCelsius}\qquad T：22、22、22、21、\qty{23}{\degreeCelsius}
\bkpans{どちらも\R{平均}{へい|きん}22。S：2\R{乗}{じょう}の\R{合計}{ごう|けい}10、\R{分散}{ぶん|さん}2、$\sigma=\sqrt{2}=1.41\cdots$。T：2\R{乗}{じょう}の\R{合計}{ごう|けい}2、\R{分散}{ぶん|さん}0.4、$\sigma=\sqrt{0.4}=0.63\cdots$。Tのほうがそろっている}

\bkproblem 10\R{日}{にち}\R{間}{かん}の\R{通学時間}{つう|がく|じ|かん}は、\R{平均}{へい|きん}30\R{分}{ぷん}、$\sigma$ 4\R{分}{ふん}でした。（1）\R{平均}{へい|きん}$\pm\sigma$ の\R{範囲}{はん|い}（2）その\R{範囲}{はん|い}に\R{入}{はい}る\R{日数}{にっ|すう}と\%を\R{求}{もと}めましょう。\par
31、24、33、28、35、30、22、34、31、32（\R{分}{ふん}）
\bkpans{（1）\nobreak 26\R{分}{ぷん}から34\R{分}{ぷん}まで（2）\nobreak 31、33、28、30、34、31、32の7\R{日}{にち}。70\%}

\bkproblem ボルトの\R{長}{なが}さは、\R{平均}{へい|きん}\qty{20.0}{mm}、$\sigma$ \qty{0.3}{mm}でした。（1）cmで\R{表}{あらわ}すと、\R{平均}{へい|きん}と $\sigma$ は？（2）どのボルトにも\R{長}{なが}さ\qty{5}{mm}のナットをつけた\R{長}{なが}さでは？
\bkpans{（1）\nobreak \qty{10}{mm/cm}でわる。\R{平均}{へい|きん}\qty{2.00}{cm}、$\sigma$ \qty{0.03}{cm}（2）\nobreak \R{平均}{へい|きん}\qty{25.0}{mm}、$\sigma$ \qty{0.3}{mm}}

\bkproblem 1\R{回}{かい}はかるときの $\sigma$ が\qty{0.8}{g}のはかりがあります。（1）16\R{回}{かい}の\R{平均}{へい|きん}のずれは？（2）\R{平均}{へい|きん}のずれを\qty{0.1}{g}にするには\R{何回}{なん|かい}？
\bkpans{（1）\nobreak $0.8\div\sqrt{16}=0.8\div4=0.2$ で\qty{0.2}{g}（2）\nobreak $0.8\div0.1=8$ なので $\sqrt{n}=8$、$n=64$ で64\R{回}{かい}}

\bkproblemgroup{チャレンジ}{3}

\bkproblem 6\R{回}{かい}はかった\R{値}{あたい}が12、8、14、11、9、12でした。（1）\R{平均}{へい|きん}と $\sigma$ を\R{求}{もと}めましょう。（2）\R{平均}{へい|きん}$\pm\sigma$ に\R{入}{はい}るのは\R{何回}{なん|かい}ですか。
\bkpans{（1）\nobreak \R{平均}{へい|きん} $66\div6=11$。\R{偏差}{へん|さ} $+1$、$-3$、$+3$、0、$-2$、$+1$、2\R{乗}{じょう}の\R{合計}{ごう|けい}24、\R{分散}{ぶん|さん}4、$\sigma=2$（2）\nobreak 9から13まで。12、11、9、12の4\R{回}{かい}}

\bkproblem \R{温度}{おん|ど}センサーの\R{値}{あたい}は、\R{平均}{へい|きん}\qty{20.0}{\degreeCelsius}、$\sigma$ \qty{0.5}{\degreeCelsius}でした。\R{華氏}{か|し}（$^{\circ}$F）は、\R{摂氏}{せっ|し}を1.8\R{倍}{ばい}して32をたした\R{温度}{おん|ど}です（3-3：$F=1.8C+32$）。\R{華氏}{か|し}で\R{表}{あらわ}すと、\R{平均}{へい|きん}と $\sigma$ はいくつですか。
\bkpans{\R{平均}{へい|きん} $20.0\times1.8+32=68.0$ で68.0\,$^{\circ}$F。$\sigma$ は1.8\R{倍}{ばい}して $0.5\times1.8=0.9$ で0.9\,$^{\circ}$F（32をたしても\R{変}{か}わらない）}

\bkproblem ふりこの\R{時間}{じ|かん}を、\R{手}{て}でストップウォッチをおしてはかると1\R{回}{かい}の $\sigma$ は0.1\R{秒}{びょう}、\R{光}{ひかり}センサーではかると0.004\R{秒}{びょう}でした。\R{手}{て}ではかった\R{平均}{へい|きん}のずれを、\R{光}{ひかり}センサー1\R{回}{かい}と\R{同}{おな}じ0.004\R{秒}{びょう}にするには、\R{手}{て}で\R{何回}{なん|かい}はかればよいですか。
\bkpans{$0.1\div0.004=25$ なので $\sqrt{n}=25$、$n=25^2=625$ で625\R{回}{かい}}

\begin{furikaeri}
\begin{checklist}
\item \R{偏差}{へん|さ}を\R{求}{もと}め、\R{合計}{ごう|けい}が0になるわけを\R{説明}{せつ|めい}できる。\bkfurref{bk:3-14:hensa}
\item \R{分散}{ぶん|さん}と $\sigma$ を\R{求}{もと}め、ばらつきをくらべられる。\bkfurref{bk:3-14:sigma}
\item \R{平均}{へい|きん}$\pm\sigma$ と $\dfrac{\sigma}{\sqrt{n}}$ の\R{意味}{い|み}を\R{説明}{せつ|めい}できる。\bkfurref{bk:3-14:nanawari}
\end{checklist}
\end{furikaeri}

\bkfukidashi{delta}{「だいたいこのくらいずれる」が $\sigma$ で\R{言}{い}えるようになったね。3-17では、このつりがねの\R{形}{かたち}がどこから\R{来}{く}るのかを\R{調}{しら}べるよ。}
